
\documentclass{article}
\usepackage{colortbl}
\usepackage{makecell}
\usepackage{multirow}
\usepackage{supertabular}

\begin{document}

\newcounter{utterance}

\centering \large Interaction Transcript for game `clean\_up', experiment `0\_easy\_7obj\_nl', episode 1 with nemotron{-}3{-}ultra.
\vspace{24pt}

{ \footnotesize  \setcounter{utterance}{1}
\setlength{\tabcolsep}{0pt}
\begin{supertabular}{c@{$\;$}|p{.15\linewidth}@{}p{.15\linewidth}p{.15\linewidth}p{.15\linewidth}p{.15\linewidth}p{.15\linewidth}}
    \# & $\;$A & \multicolumn{4}{c}{Game Master} & $\;\:$B\\
    \hline

    \theutterance \stepcounter{utterance}  
    & & \multicolumn{4}{p{0.6\linewidth}}{
        \cellcolor[rgb]{0.9,0.9,0.9}{
            \makecell[{{p{\linewidth}}}]{
                \texttt{\tiny{[P1$\langle$GM]}}
                \texttt{Ik ben uw spelleider, en u speelt een samenwerkingsspel met het volgende raster als speelbord:} \\
\\ 
\texttt{```} \\
\texttt{ 1234567} \\
\texttt{╔═══════╗} \\
\texttt{║◌◌◌◌◌T◌║ 1} \\
\texttt{║P◌◌◌◌◌◌║ 2} \\
\texttt{║◌◌AD◌┌─╢ 3} \\
\texttt{║◌◌◌┌─┤S║ 4} \\
\texttt{║O◌◌│◌├─╢ 5} \\
\texttt{╟───┼─┘◌║ 6} \\
\texttt{║◌M◌│◌◌◌║ 7} \\
\texttt{╚═══╧═══╝} \\
\\ 
\texttt{```} \\
\\ 
\texttt{* De bovenrand toont x{-}coördinaten die naar rechts oplopen, en de rechterrand y{-}coördinaten die naar beneden oplopen.} \\
\texttt{* De volgende objecten zijn willekeurig op uw raster geplaatst: 'P', 'O', 'T', 'S', 'D', 'A', 'M'.} \\
\\ 
\\ 
\texttt{De andere speler ziet een variatie van het speelbord, waarbij de objecten op verschillende willekeurige locaties zijn geplaatst. U kunt het bord van de andere speler niet zien en zij kunnen het uwe niet zien.} \\
\\ 
\\ 
\texttt{**Doel:**} \\
\\ 
\texttt{Beide spelers moeten de objecten op hun respectievelijke achtergrond zo verplaatsen dat identieke objecten op dezelfde coördinaten terechtkomen. U moet communiceren met de andere speler om een gemeenschappelijke doelconfiguratie overeen te komen.} \\
\\ 
\\ 
\texttt{**Regels:**} \\
\\ 
\texttt{* In elke beurt kunt u precies één van de volgende twee commando's sturen:} \\
\texttt{ 1. `ZEG: <BERICHT>`: om een bericht (alles tot de volgende regeleinde) naar de andere speler te sturen. Ik zal het doorsturen naar uw partner.} \\
\texttt{2. `BEWEEG: <OBJECT>, (<X>, <Y>)`: om een object naar een nieuwe positie te verplaatsen, waarbij `<X>` de kolom is en `<Y>` de rij. Ik zal u laten weten of uw zet succesvol was of niet.} \\
\texttt{* Als u zich niet aan het formaat houdt, of meerdere commando's tegelijk verstuurt, moet ik u straffen.} \\
\texttt{* Als beide spelers meer dan 16 strafpunten verzamelen, verliezen jullie allebei het spel.} \\
\texttt{* Het is van vitaal belang dat u met de andere speler communiceert over uw doelstatus! De *enige* manier waarop u uw strategie aan de andere speler kunt overbrengen is door gebruik te maken van het `ZEG: <BERICHT>` commando!} \\
\\ 
\\ 
\texttt{**Objecten verplaatsen**} \\
\\ 
\texttt{* U kunt alleen objecten verplaatsen naar cellen binnen de grenzen van het raster. De doelcel moet leeg zijn, d.w.z. deze mag alleen het symbool '◌' bevatten.} \\
\texttt{* Als u probeert een object te verplaatsen naar een plek die niet leeg is, of als u probeert het buiten het raster te verplaatsen, moet ik u straffen. U krijgt een nieuwe poging.} \\
\texttt{* Voordat u een zet doet, moet u controleren of de doelplek leeg is, en geen letter, kader of lijn bevat!} \\
\\ 
\\ 
\texttt{**Einde van het spel**} \\
\\ 
\texttt{Als u denkt dat u het doel bereikt hebt om alle objecten uit te lijnen, kunt u de andere speler vragen het spel te beëindigen door `ZEG: klaar?` te sturen. Als de andere speler u vraagt om het spel te beëindigen en u antwoordt met `ZEG: klaar!`, zal het spel eindigen.} \\
\\ 
\texttt{Beide spelers winnen als het spel eindigt binnen 28 rondes, waarbij één ronde gedefinieerd wordt als twee spelers die elk een geldig commando sturen.} \\
\\ 
\texttt{**Score:**} \\
\\ 
\texttt{Hoe dichter de identieke objecten bij elkaar liggen op beide speelborden, hoe meer punten u krijgt. Strafpunten verminderen uw punten. Kunt u het record verslaan?} \\
\\ 
\\ 
\texttt{Stuur een bericht naar de andere speler om het spel te starten!} \\
            }
        }
    }
    & \\ \\

    \theutterance \stepcounter{utterance}  
    & \multicolumn{4}{p{0.6\linewidth}}{
        \cellcolor[rgb]{0.9,0.9,0.9}{
            \makecell[{{p{\linewidth}}}]{
                \texttt{\tiny{[P1$\rangle$GM]}}
                \texttt{ZEG: Hallo! Ik ben klaar om te beginnen. Laten we onze bordconfiguraties delen om een gemeenschappelijke doelconfiguratie te vinden.} \\
\\ 
\texttt{Mijn objecten en coördinaten (X, Y):} \\
\texttt{P: (1, 2)} \\
\texttt{O: (1, 5)} \\
\texttt{T: (6, 1)} \\
\texttt{A: (3, 3)} \\
\texttt{D: (4, 3)} \\
\texttt{M: (2, 7)} \\
\texttt{S: (5, 4)} \\
\\ 
\texttt{Wat zijn jouw coördinaten?} \\
            }
        }
    }
    & & \\ \\

    \theutterance \stepcounter{utterance}  
    & & & \multicolumn{2}{p{0.3\linewidth}}{
        \cellcolor[rgb]{0.9,0.9,0.9}{
            \makecell[{{p{\linewidth}}}]{
                \texttt{\tiny{[GM$|$GM]}}
                \texttt{ZEG: Hallo! Ik ben klaar om te beginnen. Laten we onze bordconfiguraties delen om een gemeenschappelijke doelconfiguratie te vinden.} \\
\\ 
\texttt{Mijn objecten en coördinaten (X, Y):} \\
\texttt{P: (1, 2)} \\
\texttt{O: (1, 5)} \\
\texttt{T: (6, 1)} \\
\texttt{A: (3, 3)} \\
\texttt{D: (4, 3)} \\
\texttt{M: (2, 7)} \\
\texttt{S: (5, 4)} \\
\\ 
\texttt{Wat zijn jouw coördinaten?} \\
            }
        }
    }
    & & \\ \\

    \theutterance \stepcounter{utterance}  
    & & & \multicolumn{2}{p{0.3\linewidth}}{
        \cellcolor[rgb]{0.9,0.9,0.9}{
            \makecell[{{p{\linewidth}}}]{
                \texttt{\tiny{[GM$|$GM]}}
                \texttt{Invalid format: tail is not empty: '} \\
\\ 
\texttt{Mijn objecten en coördinaten (X, Y):} \\
\texttt{P: (1, 2)} \\
\texttt{O: (1, 5)} \\
\texttt{T: (6, 1)} \\
\texttt{A: (3, 3)} \\
\texttt{D: (4, 3)} \\
\texttt{M: (2, 7)} \\
\texttt{S: (5, 4)} \\
\\ 
\texttt{Wat zijn jouw coördinaten?} \\
            }
        }
    }
    & & \\ \\

    \theutterance \stepcounter{utterance}  
    & & \multicolumn{4}{p{0.6\linewidth}}{
        \cellcolor[rgb]{0.9,0.9,0.9}{
            \makecell[{{p{\linewidth}}}]{
                \texttt{\tiny{[P1$\langle$GM]}}
                \texttt{Straf: Je bericht mag niets na het commando bevatten!} \\
\\ 
\texttt{Zorg ervoor dat je antwoord alleen `ZEG: <BERICHT>` of `BEWEEG: <OBJECT>, (<X>, <Y>)` bevat, en niets anders!} \\
\\ 
\texttt{Jullie hebben gezamenlijk 1 van 16 straffen verzameld.} \\
\texttt{Probeer het opnieuw!} \\
            }
        }
    }
    & \\ \\

    \theutterance \stepcounter{utterance}  
    & \multicolumn{4}{p{0.6\linewidth}}{
        \cellcolor[rgb]{0.9,0.9,0.9}{
            \makecell[{{p{\linewidth}}}]{
                \texttt{\tiny{[P1$\rangle$GM]}}
                \texttt{ZEG: Hallo! Ik ben klaar om te beginnen. Mijn objecten en coördinaten (X, Y): P: (1, 2), O: (1, 5), T: (6, 1), A: (3, 3), D: (4, 3), M: (2, 7), S: (5, 4). Wat zijn jouw coördinaten?} \\
            }
        }
    }
    & & \\ \\

    \theutterance \stepcounter{utterance}  
    & & & \multicolumn{2}{p{0.3\linewidth}}{
        \cellcolor[rgb]{0.9,0.9,0.9}{
            \makecell[{{p{\linewidth}}}]{
                \texttt{\tiny{[GM$|$GM]}}
                \texttt{ZEG: Hallo! Ik ben klaar om te beginnen. Mijn objecten en coördinaten (X, Y): P: (1, 2), O: (1, 5), T: (6, 1), A: (3, 3), D: (4, 3), M: (2, 7), S: (5, 4). Wat zijn jouw coördinaten?} \\
            }
        }
    }
    & & \\ \\

    \theutterance \stepcounter{utterance}  
    & & \multicolumn{4}{p{0.6\linewidth}}{
        \cellcolor[rgb]{0.9,0.9,0.9}{
            \makecell[{{p{\linewidth}}}]{
                \texttt{\tiny{[GM$\rangle$P2]}}
                \texttt{Ik ben uw spelleider, en u speelt een samenwerkingsspel met het volgende raster als speelbord:} \\
\\ 
\texttt{```} \\
\texttt{ 1234567} \\
\texttt{╔═══════╗} \\
\texttt{║◌◌◌◌◌S◌║ 1} \\
\texttt{║◌◌◌◌O◌◌║ 2} \\
\texttt{║◌◌◌M◌┌─╢ 3} \\
\texttt{║◌◌A┌─┤◌║ 4} \\
\texttt{║◌◌◌│◌├─╢ 5} \\
\texttt{╟───┼─┘◌║ 6} \\
\texttt{║D◌T│P◌◌║ 7} \\
\texttt{╚═══╧═══╝} \\
\\ 
\texttt{```} \\
\\ 
\texttt{* De bovenrand toont x{-}coördinaten die naar rechts oplopen, en de rechterrand y{-}coördinaten die naar beneden oplopen.} \\
\texttt{* De volgende objecten zijn willekeurig op uw raster geplaatst: 'P', 'O', 'T', 'S', 'D', 'A', 'M'.} \\
\\ 
\\ 
\texttt{De andere speler ziet een variatie van het speelbord, waarbij de objecten op verschillende willekeurige locaties zijn geplaatst. U kunt het bord van de andere speler niet zien en zij kunnen het uwe niet zien.} \\
\\ 
\\ 
\texttt{**Doel:**} \\
\\ 
\texttt{Beide spelers moeten de objecten op hun respectievelijke achtergrond zo verplaatsen dat identieke objecten op dezelfde coördinaten terechtkomen. U moet communiceren met de andere speler om een gemeenschappelijke doelconfiguratie overeen te komen.} \\
\\ 
\\ 
\texttt{**Regels:**} \\
\\ 
\texttt{* In elke beurt kunt u precies één van de volgende twee commando's sturen:} \\
\texttt{ 1. `ZEG: <BERICHT>`: om een bericht (alles tot de volgende regeleinde) naar de andere speler te sturen. Ik zal het doorsturen naar uw partner.} \\
\texttt{2. `BEWEEG: <OBJECT>, (<X>, <Y>)`: om een object naar een nieuwe positie te verplaatsen, waarbij `<X>` de kolom is en `<Y>` de rij. Ik zal u laten weten of uw zet succesvol was of niet.} \\
\texttt{* Als u zich niet aan het formaat houdt, of meerdere commando's tegelijk verstuurt, moet ik u straffen.} \\
\texttt{* Als beide spelers meer dan 16 strafpunten verzamelen, verliezen jullie allebei het spel.} \\
\texttt{* Het is van vitaal belang dat u met de andere speler communiceert over uw doelstatus! De *enige* manier waarop u uw strategie aan de andere speler kunt overbrengen is door gebruik te maken van het `ZEG: <BERICHT>` commando!} \\
\\ 
\\ 
\texttt{**Objecten verplaatsen**} \\
\\ 
\texttt{* U kunt alleen objecten verplaatsen naar cellen binnen de grenzen van het raster. De doelcel moet leeg zijn, d.w.z. deze mag alleen het symbool '◌' bevatten.} \\
\texttt{* Als u probeert een object te verplaatsen naar een plek die niet leeg is, of als u probeert het buiten het raster te verplaatsen, moet ik u straffen. U krijgt een nieuwe poging.} \\
\texttt{* Voordat u een zet doet, moet u controleren of de doelplek leeg is, en geen letter, kader of lijn bevat!} \\
\\ 
\\ 
\texttt{**Einde van het spel**} \\
\\ 
\texttt{Als u denkt dat u het doel bereikt hebt om alle objecten uit te lijnen, kunt u de andere speler vragen het spel te beëindigen door `ZEG: klaar?` te sturen. Als de andere speler u vraagt om het spel te beëindigen en u antwoordt met `ZEG: klaar!`, zal het spel eindigen.} \\
\\ 
\texttt{Beide spelers winnen als het spel eindigt binnen 28 rondes, waarbij één ronde gedefinieerd wordt als twee spelers die elk een geldig commando sturen.} \\
\\ 
\texttt{**Score:**} \\
\\ 
\texttt{Hoe dichter de identieke objecten bij elkaar liggen op beide speelborden, hoe meer punten u krijgt. Strafpunten verminderen uw punten. Kunt u het record verslaan?} \\
\\ 
\texttt{De andere speler begon het spel door dit bericht te sturen:} \\
\\ 
\texttt{"Hallo! Ik ben klaar om te beginnen. Mijn objecten en coördinaten (X, Y): P: (1, 2), O: (1, 5), T: (6, 1), A: (3, 3), D: (4, 3), M: (2, 7), S: (5, 4). Wat zijn jouw coördinaten?"} \\
\\ 
\texttt{Wat is uw eerste opdracht?} \\
            }
        }
    }
    & \\ \\

    \theutterance \stepcounter{utterance}  
    & & & \multicolumn{4}{p{0.6\linewidth}}{
        \cellcolor[rgb]{0.9,0.9,0.9}{
            \makecell[{{p{\linewidth}}}]{
                \texttt{\tiny{[GM$\langle$P2]}}
                \texttt{ZEG: Hallo! Mijn objecten en coördinaten (X, Y): P: (6, 7), O: (5, 2), T: (4, 7), A: (3, 4), D: (1, 7), M: (4, 3), S: (6, 1). Waar willen we de objecten neerzetten? De middelste kolommen (3, 4, 5) en rijen (3, 4, 5) lijken me een goed doelgebied omdat daar ruimte is.} \\
            }
        }
    }
     \\ \\

    \theutterance \stepcounter{utterance}  
    & & & \multicolumn{2}{p{0.3\linewidth}}{
        \cellcolor[rgb]{0.9,0.9,0.9}{
            \makecell[{{p{\linewidth}}}]{
                \texttt{\tiny{[GM$|$GM]}}
                \texttt{ZEG: Hallo! Mijn objecten en coördinaten (X, Y): P: (6, 7), O: (5, 2), T: (4, 7), A: (3, 4), D: (1, 7), M: (4, 3), S: (6, 1). Waar willen we de objecten neerzetten? De middelste kolommen (3, 4, 5) en rijen (3, 4, 5) lijken me een goed doelgebied omdat daar ruimte is.} \\
            }
        }
    }
    & & \\ \\

    \theutterance \stepcounter{utterance}  
    & & \multicolumn{4}{p{0.6\linewidth}}{
        \cellcolor[rgb]{0.9,0.9,0.9}{
            \makecell[{{p{\linewidth}}}]{
                \texttt{\tiny{[P1$\langle$GM]}}
                \texttt{Je bericht is doorgestuurd naar de andere speler.} \\
\\ 
\\ 
\texttt{Je speelt momenteel ronde 1 van maximaal 28.} \\
\texttt{Jullie hebben gezamenlijk 1 van 16 straffen verzameld.} \\
\texttt{De andere speler stuurde dit bericht:} \\
\\ 
\\ 
\texttt{"Hallo! Mijn objecten en coördinaten (X, Y): P: (6, 7), O: (5, 2), T: (4, 7), A: (3, 4), D: (1, 7), M: (4, 3), S: (6, 1). Waar willen we de objecten neerzetten? De middelste kolommen (3, 4, 5) en rijen (3, 4, 5) lijken me een goed doelgebied omdat daar ruimte is."} \\
\\ 
\\ 
\texttt{Wat is je volgende opdracht?} \\
            }
        }
    }
    & \\ \\

    \theutterance \stepcounter{utterance}  
    & \multicolumn{4}{p{0.6\linewidth}}{
        \cellcolor[rgb]{0.9,0.9,0.9}{
            \makecell[{{p{\linewidth}}}]{
                \texttt{\tiny{[P1$\rangle$GM]}}
                \texttt{ZEG: Goed plan, de middelste cellen (kolommen 3{-}5, rijen 3{-}5) zijn ideaal. Laten we de volgende doelcoördinaten afspreken: P:(3,3), O:(4,3), T:(5,3), A:(3,4), D:(4,4), M:(5,4), S:(4,5). Als je akkoord bent, begin ik met verplaatsen. Akkoord?} \\
            }
        }
    }
    & & \\ \\

    \theutterance \stepcounter{utterance}  
    & & & \multicolumn{2}{p{0.3\linewidth}}{
        \cellcolor[rgb]{0.9,0.9,0.9}{
            \makecell[{{p{\linewidth}}}]{
                \texttt{\tiny{[GM$|$GM]}}
                \texttt{ZEG: Goed plan, de middelste cellen (kolommen 3{-}5, rijen 3{-}5) zijn ideaal. Laten we de volgende doelcoördinaten afspreken: P:(3,3), O:(4,3), T:(5,3), A:(3,4), D:(4,4), M:(5,4), S:(4,5). Als je akkoord bent, begin ik met verplaatsen. Akkoord?} \\
            }
        }
    }
    & & \\ \\

    \theutterance \stepcounter{utterance}  
    & & \multicolumn{4}{p{0.6\linewidth}}{
        \cellcolor[rgb]{0.9,0.9,0.9}{
            \makecell[{{p{\linewidth}}}]{
                \texttt{\tiny{[GM$\rangle$P2]}}
                \texttt{Je bericht is doorgestuurd naar de andere speler.} \\
\\ 
\\ 
\texttt{Je speelt momenteel ronde 2 van maximaal 28.} \\
\texttt{Jullie hebben gezamenlijk 1 van 16 straffen verzameld.} \\
\texttt{De andere speler stuurde dit bericht:} \\
\\ 
\\ 
\texttt{"Goed plan, de middelste cellen (kolommen 3{-}5, rijen 3{-}5) zijn ideaal. Laten we de volgende doelcoördinaten afspreken: P:(3,3), O:(4,3), T:(5,3), A:(3,4), D:(4,4), M:(5,4), S:(4,5). Als je akkoord bent, begin ik met verplaatsen. Akkoord?"} \\
\\ 
\\ 
\texttt{Wat is je volgende opdracht?} \\
            }
        }
    }
    & \\ \\

    \theutterance \stepcounter{utterance}  
    & & & \multicolumn{4}{p{0.6\linewidth}}{
        \cellcolor[rgb]{0.9,0.9,0.9}{
            \makecell[{{p{\linewidth}}}]{
                \texttt{\tiny{[GM$\langle$P2]}}
                \texttt{ZEG: Akkoord met de doelcoördinaten: P:(3,3), O:(4,3), T:(5,3), A:(3,4), D:(4,4), M:(5,4), S:(4,5). Ik begin met verplaatsen.} \\
            }
        }
    }
     \\ \\

    \theutterance \stepcounter{utterance}  
    & & & \multicolumn{2}{p{0.3\linewidth}}{
        \cellcolor[rgb]{0.9,0.9,0.9}{
            \makecell[{{p{\linewidth}}}]{
                \texttt{\tiny{[GM$|$GM]}}
                \texttt{ZEG: Akkoord met de doelcoördinaten: P:(3,3), O:(4,3), T:(5,3), A:(3,4), D:(4,4), M:(5,4), S:(4,5). Ik begin met verplaatsen.} \\
            }
        }
    }
    & & \\ \\

    \theutterance \stepcounter{utterance}  
    & & \multicolumn{4}{p{0.6\linewidth}}{
        \cellcolor[rgb]{0.9,0.9,0.9}{
            \makecell[{{p{\linewidth}}}]{
                \texttt{\tiny{[P1$\langle$GM]}}
                \texttt{Je bericht is doorgestuurd naar de andere speler.} \\
\\ 
\\ 
\texttt{Je speelt momenteel ronde 2 van maximaal 28.} \\
\texttt{Jullie hebben gezamenlijk 1 van 16 straffen verzameld.} \\
\texttt{De andere speler stuurde dit bericht:} \\
\\ 
\\ 
\texttt{"Akkoord met de doelcoördinaten: P:(3,3), O:(4,3), T:(5,3), A:(3,4), D:(4,4), M:(5,4), S:(4,5). Ik begin met verplaatsen."} \\
\\ 
\\ 
\texttt{Wat is je volgende opdracht?} \\
            }
        }
    }
    & \\ \\

    \theutterance \stepcounter{utterance}  
    & \multicolumn{4}{p{0.6\linewidth}}{
        \cellcolor[rgb]{0.9,0.9,0.9}{
            \makecell[{{p{\linewidth}}}]{
                \texttt{\tiny{[P1$\rangle$GM]}}
                \texttt{BEWEEG: P, (3, 3)} \\
            }
        }
    }
    & & \\ \\

    \theutterance \stepcounter{utterance}  
    & & & \multicolumn{2}{p{0.3\linewidth}}{
        \cellcolor[rgb]{0.9,0.9,0.9}{
            \makecell[{{p{\linewidth}}}]{
                \texttt{\tiny{[GM$|$GM]}}
                \texttt{BEWEEG: P, (3, 3)} \\
            }
        }
    }
    & & \\ \\

    \theutterance \stepcounter{utterance}  
    & & & \multicolumn{2}{p{0.3\linewidth}}{
        \cellcolor[rgb]{0.9,0.9,0.9}{
            \makecell[{{p{\linewidth}}}]{
                \texttt{\tiny{[GM$|$GM]}}
                \texttt{Straf: (3,3) is niet leeg, maar bevat 'A'.} \\
\\ 
\texttt{Jullie hebben gezamenlijk 2 van 16 straffen verzameld.} \\
\texttt{Probeer het opnieuw!} \\
            }
        }
    }
    & & \\ \\

    \theutterance \stepcounter{utterance}  
    & & \multicolumn{4}{p{0.6\linewidth}}{
        \cellcolor[rgb]{0.9,0.9,0.9}{
            \makecell[{{p{\linewidth}}}]{
                \texttt{\tiny{[P1$\langle$GM]}}
                \texttt{Straf: (3,3) is niet leeg, maar bevat 'A'.} \\
\\ 
\texttt{Jullie hebben gezamenlijk 2 van 16 straffen verzameld.} \\
\texttt{Probeer het opnieuw!} \\
            }
        }
    }
    & \\ \\

    \theutterance \stepcounter{utterance}  
    & \multicolumn{4}{p{0.6\linewidth}}{
        \cellcolor[rgb]{0.9,0.9,0.9}{
            \makecell[{{p{\linewidth}}}]{
                \texttt{\tiny{[P1$\rangle$GM]}}
                \texttt{BEWEEG: A, (3, 4)} \\
            }
        }
    }
    & & \\ \\

    \theutterance \stepcounter{utterance}  
    & & & \multicolumn{2}{p{0.3\linewidth}}{
        \cellcolor[rgb]{0.9,0.9,0.9}{
            \makecell[{{p{\linewidth}}}]{
                \texttt{\tiny{[GM$|$GM]}}
                \texttt{BEWEEG: A, (3, 4)} \\
            }
        }
    }
    & & \\ \\

    \theutterance \stepcounter{utterance}  
    & & & \multicolumn{2}{p{0.3\linewidth}}{
        \cellcolor[rgb]{0.9,0.9,0.9}{
            \makecell[{{p{\linewidth}}}]{
                \texttt{\tiny{[GM$|$GM]}}
                \texttt{A moved to (3, 4)} \\
            }
        }
    }
    & & \\ \\

    \theutterance \stepcounter{utterance}  
    & & \multicolumn{4}{p{0.6\linewidth}}{
        \cellcolor[rgb]{0.9,0.9,0.9}{
            \makecell[{{p{\linewidth}}}]{
                \texttt{\tiny{[GM$\rangle$P2]}}
                \texttt{Je bericht is doorgestuurd naar de andere speler.} \\
\\ 
\\ 
\texttt{Je speelt momenteel ronde 3 van maximaal 28.} \\
\texttt{Jullie hebben gezamenlijk 2 van 16 straffen verzameld.} \\
\texttt{De andere speler heeft een voorwerp op zijn rooster verplaatst.} \\
\\ 
\\ 
\texttt{Wat is je volgende opdracht?} \\
            }
        }
    }
    & \\ \\

    \theutterance \stepcounter{utterance}  
    & & & \multicolumn{4}{p{0.6\linewidth}}{
        \cellcolor[rgb]{0.9,0.9,0.9}{
            \makecell[{{p{\linewidth}}}]{
                \texttt{\tiny{[GM$\langle$P2]}}
                \texttt{BEWEEG: S, (4, 5)} \\
            }
        }
    }
     \\ \\

    \theutterance \stepcounter{utterance}  
    & & & \multicolumn{2}{p{0.3\linewidth}}{
        \cellcolor[rgb]{0.9,0.9,0.9}{
            \makecell[{{p{\linewidth}}}]{
                \texttt{\tiny{[GM$|$GM]}}
                \texttt{BEWEEG: S, (4, 5)} \\
            }
        }
    }
    & & \\ \\

    \theutterance \stepcounter{utterance}  
    & & & \multicolumn{2}{p{0.3\linewidth}}{
        \cellcolor[rgb]{0.9,0.9,0.9}{
            \makecell[{{p{\linewidth}}}]{
                \texttt{\tiny{[GM$|$GM]}}
                \texttt{Straf: (4,5) is niet leeg, maar bevat '│'.} \\
\\ 
\texttt{Jullie hebben gezamenlijk 3 van 16 straffen verzameld.} \\
\texttt{Probeer het opnieuw!} \\
            }
        }
    }
    & & \\ \\

    \theutterance \stepcounter{utterance}  
    & & \multicolumn{4}{p{0.6\linewidth}}{
        \cellcolor[rgb]{0.9,0.9,0.9}{
            \makecell[{{p{\linewidth}}}]{
                \texttt{\tiny{[GM$\rangle$P2]}}
                \texttt{Straf: (4,5) is niet leeg, maar bevat '│'.} \\
\\ 
\texttt{Jullie hebben gezamenlijk 3 van 16 straffen verzameld.} \\
\texttt{Probeer het opnieuw!} \\
            }
        }
    }
    & \\ \\

    \theutterance \stepcounter{utterance}  
    & & & \multicolumn{4}{p{0.6\linewidth}}{
        \cellcolor[rgb]{0.9,0.9,0.9}{
            \makecell[{{p{\linewidth}}}]{
                \texttt{\tiny{[GM$\langle$P2]}}
                \texttt{BEWEEG: O, (4, 3)} \\
            }
        }
    }
     \\ \\

    \theutterance \stepcounter{utterance}  
    & & & \multicolumn{2}{p{0.3\linewidth}}{
        \cellcolor[rgb]{0.9,0.9,0.9}{
            \makecell[{{p{\linewidth}}}]{
                \texttt{\tiny{[GM$|$GM]}}
                \texttt{BEWEEG: O, (4, 3)} \\
            }
        }
    }
    & & \\ \\

    \theutterance \stepcounter{utterance}  
    & & & \multicolumn{2}{p{0.3\linewidth}}{
        \cellcolor[rgb]{0.9,0.9,0.9}{
            \makecell[{{p{\linewidth}}}]{
                \texttt{\tiny{[GM$|$GM]}}
                \texttt{Straf: (4,3) is niet leeg, maar bevat 'M'.} \\
\\ 
\texttt{Jullie hebben gezamenlijk 4 van 16 straffen verzameld.} \\
\texttt{Probeer het opnieuw!} \\
            }
        }
    }
    & & \\ \\

    \theutterance \stepcounter{utterance}  
    & & \multicolumn{4}{p{0.6\linewidth}}{
        \cellcolor[rgb]{0.9,0.9,0.9}{
            \makecell[{{p{\linewidth}}}]{
                \texttt{\tiny{[GM$\rangle$P2]}}
                \texttt{Straf: (4,3) is niet leeg, maar bevat 'M'.} \\
\\ 
\texttt{Jullie hebben gezamenlijk 4 van 16 straffen verzameld.} \\
\texttt{Probeer het opnieuw!} \\
            }
        }
    }
    & \\ \\

    \theutterance \stepcounter{utterance}  
    & & & \multicolumn{4}{p{0.6\linewidth}}{
        \cellcolor[rgb]{0.9,0.9,0.9}{
            \makecell[{{p{\linewidth}}}]{
                \texttt{\tiny{[GM$\langle$P2]}}
                \texttt{BEWEEG: M, (5, 4)} \\
            }
        }
    }
     \\ \\

    \theutterance \stepcounter{utterance}  
    & & & \multicolumn{2}{p{0.3\linewidth}}{
        \cellcolor[rgb]{0.9,0.9,0.9}{
            \makecell[{{p{\linewidth}}}]{
                \texttt{\tiny{[GM$|$GM]}}
                \texttt{BEWEEG: M, (5, 4)} \\
            }
        }
    }
    & & \\ \\

    \theutterance \stepcounter{utterance}  
    & & & \multicolumn{2}{p{0.3\linewidth}}{
        \cellcolor[rgb]{0.9,0.9,0.9}{
            \makecell[{{p{\linewidth}}}]{
                \texttt{\tiny{[GM$|$GM]}}
                \texttt{Straf: (5,4) is niet leeg, maar bevat '─'.} \\
\\ 
\texttt{Jullie hebben gezamenlijk 5 van 16 straffen verzameld.} \\
\texttt{Probeer het opnieuw!} \\
            }
        }
    }
    & & \\ \\

    \theutterance \stepcounter{utterance}  
    & & \multicolumn{4}{p{0.6\linewidth}}{
        \cellcolor[rgb]{0.9,0.9,0.9}{
            \makecell[{{p{\linewidth}}}]{
                \texttt{\tiny{[GM$\rangle$P2]}}
                \texttt{Straf: (5,4) is niet leeg, maar bevat '─'.} \\
\\ 
\texttt{Jullie hebben gezamenlijk 5 van 16 straffen verzameld.} \\
\texttt{Probeer het opnieuw!} \\
            }
        }
    }
    & \\ \\

    \theutterance \stepcounter{utterance}  
    & & & \multicolumn{4}{p{0.6\linewidth}}{
        \cellcolor[rgb]{0.9,0.9,0.9}{
            \makecell[{{p{\linewidth}}}]{
                \texttt{\tiny{[GM$\langle$P2]}}
                \texttt{BEWEEG: P, (3, 3)} \\
            }
        }
    }
     \\ \\

    \theutterance \stepcounter{utterance}  
    & & & \multicolumn{2}{p{0.3\linewidth}}{
        \cellcolor[rgb]{0.9,0.9,0.9}{
            \makecell[{{p{\linewidth}}}]{
                \texttt{\tiny{[GM$|$GM]}}
                \texttt{BEWEEG: P, (3, 3)} \\
            }
        }
    }
    & & \\ \\

    \theutterance \stepcounter{utterance}  
    & & & \multicolumn{2}{p{0.3\linewidth}}{
        \cellcolor[rgb]{0.9,0.9,0.9}{
            \makecell[{{p{\linewidth}}}]{
                \texttt{\tiny{[GM$|$GM]}}
                \texttt{P moved to (3, 3)} \\
            }
        }
    }
    & & \\ \\

    \theutterance \stepcounter{utterance}  
    & & \multicolumn{4}{p{0.6\linewidth}}{
        \cellcolor[rgb]{0.9,0.9,0.9}{
            \makecell[{{p{\linewidth}}}]{
                \texttt{\tiny{[P1$\langle$GM]}}
                \texttt{Verplaatst 'A' naar (3,4) met succes. Uw bijgewerkte raster ziet er als volgt uit:} \\
\texttt{```} \\
\texttt{ 1234567} \\
\texttt{╔═══════╗} \\
\texttt{║◌◌◌◌◌T◌║ 1} \\
\texttt{║P◌◌◌◌◌◌║ 2} \\
\texttt{║◌◌◌D◌┌─╢ 3} \\
\texttt{║◌◌A┌─┤S║ 4} \\
\texttt{║O◌◌│◌├─╢ 5} \\
\texttt{╟───┼─┘◌║ 6} \\
\texttt{║◌M◌│◌◌◌║ 7} \\
\texttt{╚═══╧═══╝} \\
\texttt{```} \\
\texttt{Je speelt momenteel ronde 3 van maximaal 28.} \\
\texttt{Jullie hebben gezamenlijk 5 van 16 straffen verzameld.} \\
\texttt{De andere speler heeft een voorwerp op zijn rooster verplaatst.} \\
\\ 
\\ 
\texttt{Wat is je volgende opdracht?} \\
            }
        }
    }
    & \\ \\

    \theutterance \stepcounter{utterance}  
    & \multicolumn{4}{p{0.6\linewidth}}{
        \cellcolor[rgb]{0.9,0.9,0.9}{
            \makecell[{{p{\linewidth}}}]{
                \texttt{\tiny{[P1$\rangle$GM]}}
                \texttt{BEWEEG: P, (3, 3)} \\
            }
        }
    }
    & & \\ \\

    \theutterance \stepcounter{utterance}  
    & & & \multicolumn{2}{p{0.3\linewidth}}{
        \cellcolor[rgb]{0.9,0.9,0.9}{
            \makecell[{{p{\linewidth}}}]{
                \texttt{\tiny{[GM$|$GM]}}
                \texttt{BEWEEG: P, (3, 3)} \\
            }
        }
    }
    & & \\ \\

    \theutterance \stepcounter{utterance}  
    & & & \multicolumn{2}{p{0.3\linewidth}}{
        \cellcolor[rgb]{0.9,0.9,0.9}{
            \makecell[{{p{\linewidth}}}]{
                \texttt{\tiny{[GM$|$GM]}}
                \texttt{P moved to (3, 3)} \\
            }
        }
    }
    & & \\ \\

    \theutterance \stepcounter{utterance}  
    & & \multicolumn{4}{p{0.6\linewidth}}{
        \cellcolor[rgb]{0.9,0.9,0.9}{
            \makecell[{{p{\linewidth}}}]{
                \texttt{\tiny{[GM$\rangle$P2]}}
                \texttt{Verplaatst 'P' naar (3,3) met succes. Uw bijgewerkte raster ziet er als volgt uit:} \\
\texttt{```} \\
\texttt{ 1234567} \\
\texttt{╔═══════╗} \\
\texttt{║◌◌◌◌◌S◌║ 1} \\
\texttt{║◌◌◌◌O◌◌║ 2} \\
\texttt{║◌◌PM◌┌─╢ 3} \\
\texttt{║◌◌A┌─┤◌║ 4} \\
\texttt{║◌◌◌│◌├─╢ 5} \\
\texttt{╟───┼─┘◌║ 6} \\
\texttt{║D◌T│◌◌◌║ 7} \\
\texttt{╚═══╧═══╝} \\
\texttt{```} \\
\texttt{Je speelt momenteel ronde 4 van maximaal 28.} \\
\texttt{Jullie hebben gezamenlijk 5 van 16 straffen verzameld.} \\
\texttt{De andere speler heeft een voorwerp op zijn rooster verplaatst.} \\
\\ 
\\ 
\texttt{Wat is je volgende opdracht?} \\
            }
        }
    }
    & \\ \\

    \theutterance \stepcounter{utterance}  
    & & & \multicolumn{4}{p{0.6\linewidth}}{
        \cellcolor[rgb]{0.9,0.9,0.9}{
            \makecell[{{p{\linewidth}}}]{
                \texttt{\tiny{[GM$\langle$P2]}}
                \texttt{BEWEEG: D, (4, 4)} \\
            }
        }
    }
     \\ \\

    \theutterance \stepcounter{utterance}  
    & & & \multicolumn{2}{p{0.3\linewidth}}{
        \cellcolor[rgb]{0.9,0.9,0.9}{
            \makecell[{{p{\linewidth}}}]{
                \texttt{\tiny{[GM$|$GM]}}
                \texttt{BEWEEG: D, (4, 4)} \\
            }
        }
    }
    & & \\ \\

    \theutterance \stepcounter{utterance}  
    & & & \multicolumn{2}{p{0.3\linewidth}}{
        \cellcolor[rgb]{0.9,0.9,0.9}{
            \makecell[{{p{\linewidth}}}]{
                \texttt{\tiny{[GM$|$GM]}}
                \texttt{Straf: (4,4) is niet leeg, maar bevat '┌'.} \\
\\ 
\texttt{Jullie hebben gezamenlijk 6 van 16 straffen verzameld.} \\
\texttt{Probeer het opnieuw!} \\
            }
        }
    }
    & & \\ \\

    \theutterance \stepcounter{utterance}  
    & & \multicolumn{4}{p{0.6\linewidth}}{
        \cellcolor[rgb]{0.9,0.9,0.9}{
            \makecell[{{p{\linewidth}}}]{
                \texttt{\tiny{[GM$\rangle$P2]}}
                \texttt{Straf: (4,4) is niet leeg, maar bevat '┌'.} \\
\\ 
\texttt{Jullie hebben gezamenlijk 6 van 16 straffen verzameld.} \\
\texttt{Probeer het opnieuw!} \\
            }
        }
    }
    & \\ \\

    \theutterance \stepcounter{utterance}  
    & & & \multicolumn{4}{p{0.6\linewidth}}{
        \cellcolor[rgb]{0.9,0.9,0.9}{
            \makecell[{{p{\linewidth}}}]{
                \texttt{\tiny{[GM$\langle$P2]}}
                \texttt{BEWEEG: A, (3, 4)} \\
            }
        }
    }
     \\ \\

    \theutterance \stepcounter{utterance}  
    & & & \multicolumn{2}{p{0.3\linewidth}}{
        \cellcolor[rgb]{0.9,0.9,0.9}{
            \makecell[{{p{\linewidth}}}]{
                \texttt{\tiny{[GM$|$GM]}}
                \texttt{BEWEEG: A, (3, 4)} \\
            }
        }
    }
    & & \\ \\

    \theutterance \stepcounter{utterance}  
    & & & \multicolumn{2}{p{0.3\linewidth}}{
        \cellcolor[rgb]{0.9,0.9,0.9}{
            \makecell[{{p{\linewidth}}}]{
                \texttt{\tiny{[GM$|$GM]}}
                \texttt{Straf: (3,4) is niet leeg, maar bevat 'A'.} \\
\\ 
\texttt{Jullie hebben gezamenlijk 7 van 16 straffen verzameld.} \\
\texttt{Probeer het opnieuw!} \\
            }
        }
    }
    & & \\ \\

    \theutterance \stepcounter{utterance}  
    & & \multicolumn{4}{p{0.6\linewidth}}{
        \cellcolor[rgb]{0.9,0.9,0.9}{
            \makecell[{{p{\linewidth}}}]{
                \texttt{\tiny{[GM$\rangle$P2]}}
                \texttt{Straf: (3,4) is niet leeg, maar bevat 'A'.} \\
\\ 
\texttt{Jullie hebben gezamenlijk 7 van 16 straffen verzameld.} \\
\texttt{Probeer het opnieuw!} \\
            }
        }
    }
    & \\ \\

    \theutterance \stepcounter{utterance}  
    & & & \multicolumn{4}{p{0.6\linewidth}}{
        \cellcolor[rgb]{0.9,0.9,0.9}{
            \makecell[{{p{\linewidth}}}]{
                \texttt{\tiny{[GM$\langle$P2]}}
                \texttt{BEWEEG: T, (5, 3)} \\
            }
        }
    }
     \\ \\

    \theutterance \stepcounter{utterance}  
    & & & \multicolumn{2}{p{0.3\linewidth}}{
        \cellcolor[rgb]{0.9,0.9,0.9}{
            \makecell[{{p{\linewidth}}}]{
                \texttt{\tiny{[GM$|$GM]}}
                \texttt{BEWEEG: T, (5, 3)} \\
            }
        }
    }
    & & \\ \\

    \theutterance \stepcounter{utterance}  
    & & & \multicolumn{2}{p{0.3\linewidth}}{
        \cellcolor[rgb]{0.9,0.9,0.9}{
            \makecell[{{p{\linewidth}}}]{
                \texttt{\tiny{[GM$|$GM]}}
                \texttt{T moved to (5, 3)} \\
            }
        }
    }
    & & \\ \\

    \theutterance \stepcounter{utterance}  
    & & \multicolumn{4}{p{0.6\linewidth}}{
        \cellcolor[rgb]{0.9,0.9,0.9}{
            \makecell[{{p{\linewidth}}}]{
                \texttt{\tiny{[P1$\langle$GM]}}
                \texttt{Verplaatst 'P' naar (3,3) met succes. Uw bijgewerkte raster ziet er als volgt uit:} \\
\texttt{```} \\
\texttt{ 1234567} \\
\texttt{╔═══════╗} \\
\texttt{║◌◌◌◌◌T◌║ 1} \\
\texttt{║◌◌◌◌◌◌◌║ 2} \\
\texttt{║◌◌PD◌┌─╢ 3} \\
\texttt{║◌◌A┌─┤S║ 4} \\
\texttt{║O◌◌│◌├─╢ 5} \\
\texttt{╟───┼─┘◌║ 6} \\
\texttt{║◌M◌│◌◌◌║ 7} \\
\texttt{╚═══╧═══╝} \\
\texttt{```} \\
\texttt{Je speelt momenteel ronde 4 van maximaal 28.} \\
\texttt{Jullie hebben gezamenlijk 7 van 16 straffen verzameld.} \\
\texttt{De andere speler heeft een voorwerp op zijn rooster verplaatst.} \\
\\ 
\\ 
\texttt{Wat is je volgende opdracht?} \\
            }
        }
    }
    & \\ \\

    \theutterance \stepcounter{utterance}  
    & \multicolumn{4}{p{0.6\linewidth}}{
        \cellcolor[rgb]{0.9,0.9,0.9}{
            \makecell[{{p{\linewidth}}}]{
                \texttt{\tiny{[P1$\rangle$GM]}}
                \texttt{BEWEEG: D, (4, 4)} \\
            }
        }
    }
    & & \\ \\

    \theutterance \stepcounter{utterance}  
    & & & \multicolumn{2}{p{0.3\linewidth}}{
        \cellcolor[rgb]{0.9,0.9,0.9}{
            \makecell[{{p{\linewidth}}}]{
                \texttt{\tiny{[GM$|$GM]}}
                \texttt{BEWEEG: D, (4, 4)} \\
            }
        }
    }
    & & \\ \\

    \theutterance \stepcounter{utterance}  
    & & & \multicolumn{2}{p{0.3\linewidth}}{
        \cellcolor[rgb]{0.9,0.9,0.9}{
            \makecell[{{p{\linewidth}}}]{
                \texttt{\tiny{[GM$|$GM]}}
                \texttt{Straf: (4,4) is niet leeg, maar bevat '┌'.} \\
\\ 
\texttt{Jullie hebben gezamenlijk 8 van 16 straffen verzameld.} \\
\texttt{Probeer het opnieuw!} \\
            }
        }
    }
    & & \\ \\

    \theutterance \stepcounter{utterance}  
    & & \multicolumn{4}{p{0.6\linewidth}}{
        \cellcolor[rgb]{0.9,0.9,0.9}{
            \makecell[{{p{\linewidth}}}]{
                \texttt{\tiny{[P1$\langle$GM]}}
                \texttt{Straf: (4,4) is niet leeg, maar bevat '┌'.} \\
\\ 
\texttt{Jullie hebben gezamenlijk 8 van 16 straffen verzameld.} \\
\texttt{Probeer het opnieuw!} \\
            }
        }
    }
    & \\ \\

    \theutterance \stepcounter{utterance}  
    & \multicolumn{4}{p{0.6\linewidth}}{
        \cellcolor[rgb]{0.9,0.9,0.9}{
            \makecell[{{p{\linewidth}}}]{
                \texttt{\tiny{[P1$\rangle$GM]}}
                \texttt{BEWEEG: D, (4, 3)} \\
            }
        }
    }
    & & \\ \\

    \theutterance \stepcounter{utterance}  
    & & & \multicolumn{2}{p{0.3\linewidth}}{
        \cellcolor[rgb]{0.9,0.9,0.9}{
            \makecell[{{p{\linewidth}}}]{
                \texttt{\tiny{[GM$|$GM]}}
                \texttt{BEWEEG: D, (4, 3)} \\
            }
        }
    }
    & & \\ \\

    \theutterance \stepcounter{utterance}  
    & & & \multicolumn{2}{p{0.3\linewidth}}{
        \cellcolor[rgb]{0.9,0.9,0.9}{
            \makecell[{{p{\linewidth}}}]{
                \texttt{\tiny{[GM$|$GM]}}
                \texttt{Straf: (4,3) is niet leeg, maar bevat 'D'.} \\
\\ 
\texttt{Jullie hebben gezamenlijk 9 van 16 straffen verzameld.} \\
\texttt{Probeer het opnieuw!} \\
            }
        }
    }
    & & \\ \\

    \theutterance \stepcounter{utterance}  
    & & \multicolumn{4}{p{0.6\linewidth}}{
        \cellcolor[rgb]{0.9,0.9,0.9}{
            \makecell[{{p{\linewidth}}}]{
                \texttt{\tiny{[P1$\langle$GM]}}
                \texttt{Straf: (4,3) is niet leeg, maar bevat 'D'.} \\
\\ 
\texttt{Jullie hebben gezamenlijk 9 van 16 straffen verzameld.} \\
\texttt{Probeer het opnieuw!} \\
            }
        }
    }
    & \\ \\

    \theutterance \stepcounter{utterance}  
    & \multicolumn{4}{p{0.6\linewidth}}{
        \cellcolor[rgb]{0.9,0.9,0.9}{
            \makecell[{{p{\linewidth}}}]{
                \texttt{\tiny{[P1$\rangle$GM]}}
                \texttt{BEWEEG: S, (4, 5)} \\
            }
        }
    }
    & & \\ \\

    \theutterance \stepcounter{utterance}  
    & & & \multicolumn{2}{p{0.3\linewidth}}{
        \cellcolor[rgb]{0.9,0.9,0.9}{
            \makecell[{{p{\linewidth}}}]{
                \texttt{\tiny{[GM$|$GM]}}
                \texttt{BEWEEG: S, (4, 5)} \\
            }
        }
    }
    & & \\ \\

    \theutterance \stepcounter{utterance}  
    & & & \multicolumn{2}{p{0.3\linewidth}}{
        \cellcolor[rgb]{0.9,0.9,0.9}{
            \makecell[{{p{\linewidth}}}]{
                \texttt{\tiny{[GM$|$GM]}}
                \texttt{Straf: (4,5) is niet leeg, maar bevat '│'.} \\
\\ 
\texttt{Jullie hebben gezamenlijk 10 van 16 straffen verzameld.} \\
\texttt{Probeer het opnieuw!} \\
            }
        }
    }
    & & \\ \\

    \theutterance \stepcounter{utterance}  
    & & \multicolumn{4}{p{0.6\linewidth}}{
        \cellcolor[rgb]{0.9,0.9,0.9}{
            \makecell[{{p{\linewidth}}}]{
                \texttt{\tiny{[P1$\langle$GM]}}
                \texttt{Straf: (4,5) is niet leeg, maar bevat '│'.} \\
\\ 
\texttt{Jullie hebben gezamenlijk 10 van 16 straffen verzameld.} \\
\texttt{Probeer het opnieuw!} \\
            }
        }
    }
    & \\ \\

    \theutterance \stepcounter{utterance}  
    & \multicolumn{4}{p{0.6\linewidth}}{
        \cellcolor[rgb]{0.9,0.9,0.9}{
            \makecell[{{p{\linewidth}}}]{
                \texttt{\tiny{[P1$\rangle$GM]}}
                \texttt{BEWEEG: O, (4, 3)} \\
            }
        }
    }
    & & \\ \\

    \theutterance \stepcounter{utterance}  
    & & & \multicolumn{2}{p{0.3\linewidth}}{
        \cellcolor[rgb]{0.9,0.9,0.9}{
            \makecell[{{p{\linewidth}}}]{
                \texttt{\tiny{[GM$|$GM]}}
                \texttt{BEWEEG: O, (4, 3)} \\
            }
        }
    }
    & & \\ \\

    \theutterance \stepcounter{utterance}  
    & & & \multicolumn{2}{p{0.3\linewidth}}{
        \cellcolor[rgb]{0.9,0.9,0.9}{
            \makecell[{{p{\linewidth}}}]{
                \texttt{\tiny{[GM$|$GM]}}
                \texttt{Straf: (4,3) is niet leeg, maar bevat 'D'.} \\
\\ 
\texttt{Jullie hebben gezamenlijk 11 van 16 straffen verzameld.} \\
\texttt{Probeer het opnieuw!} \\
            }
        }
    }
    & & \\ \\

    \theutterance \stepcounter{utterance}  
    & & \multicolumn{4}{p{0.6\linewidth}}{
        \cellcolor[rgb]{0.9,0.9,0.9}{
            \makecell[{{p{\linewidth}}}]{
                \texttt{\tiny{[P1$\langle$GM]}}
                \texttt{Straf: (4,3) is niet leeg, maar bevat 'D'.} \\
\\ 
\texttt{Jullie hebben gezamenlijk 11 van 16 straffen verzameld.} \\
\texttt{Probeer het opnieuw!} \\
            }
        }
    }
    & \\ \\

    \theutterance \stepcounter{utterance}  
    & \multicolumn{4}{p{0.6\linewidth}}{
        \cellcolor[rgb]{0.9,0.9,0.9}{
            \makecell[{{p{\linewidth}}}]{
                \texttt{\tiny{[P1$\rangle$GM]}}
                \texttt{BEWEEG: T, (5, 3)} \\
            }
        }
    }
    & & \\ \\

    \theutterance \stepcounter{utterance}  
    & & & \multicolumn{2}{p{0.3\linewidth}}{
        \cellcolor[rgb]{0.9,0.9,0.9}{
            \makecell[{{p{\linewidth}}}]{
                \texttt{\tiny{[GM$|$GM]}}
                \texttt{BEWEEG: T, (5, 3)} \\
            }
        }
    }
    & & \\ \\

    \theutterance \stepcounter{utterance}  
    & & & \multicolumn{2}{p{0.3\linewidth}}{
        \cellcolor[rgb]{0.9,0.9,0.9}{
            \makecell[{{p{\linewidth}}}]{
                \texttt{\tiny{[GM$|$GM]}}
                \texttt{T moved to (5, 3)} \\
            }
        }
    }
    & & \\ \\

    \theutterance \stepcounter{utterance}  
    & & \multicolumn{4}{p{0.6\linewidth}}{
        \cellcolor[rgb]{0.9,0.9,0.9}{
            \makecell[{{p{\linewidth}}}]{
                \texttt{\tiny{[GM$\rangle$P2]}}
                \texttt{Verplaatst 'T' naar (5,3) met succes. Uw bijgewerkte raster ziet er als volgt uit:} \\
\texttt{```} \\
\texttt{ 1234567} \\
\texttt{╔═══════╗} \\
\texttt{║◌◌◌◌◌S◌║ 1} \\
\texttt{║◌◌◌◌O◌◌║ 2} \\
\texttt{║◌◌PMT┌─╢ 3} \\
\texttt{║◌◌A┌─┤◌║ 4} \\
\texttt{║◌◌◌│◌├─╢ 5} \\
\texttt{╟───┼─┘◌║ 6} \\
\texttt{║D◌◌│◌◌◌║ 7} \\
\texttt{╚═══╧═══╝} \\
\texttt{```} \\
\texttt{Je speelt momenteel ronde 5 van maximaal 28.} \\
\texttt{Jullie hebben gezamenlijk 11 van 16 straffen verzameld.} \\
\texttt{De andere speler heeft een voorwerp op zijn rooster verplaatst.} \\
\\ 
\\ 
\texttt{Wat is je volgende opdracht?} \\
            }
        }
    }
    & \\ \\

    \theutterance \stepcounter{utterance}  
    & & & \multicolumn{4}{p{0.6\linewidth}}{
        \cellcolor[rgb]{0.9,0.9,0.9}{
            \makecell[{{p{\linewidth}}}]{
                \texttt{\tiny{[GM$\langle$P2]}}
                \texttt{BEWEEG: D, (5, 4)} \\
            }
        }
    }
     \\ \\

    \theutterance \stepcounter{utterance}  
    & & & \multicolumn{2}{p{0.3\linewidth}}{
        \cellcolor[rgb]{0.9,0.9,0.9}{
            \makecell[{{p{\linewidth}}}]{
                \texttt{\tiny{[GM$|$GM]}}
                \texttt{BEWEEG: D, (5, 4)} \\
            }
        }
    }
    & & \\ \\

    \theutterance \stepcounter{utterance}  
    & & & \multicolumn{2}{p{0.3\linewidth}}{
        \cellcolor[rgb]{0.9,0.9,0.9}{
            \makecell[{{p{\linewidth}}}]{
                \texttt{\tiny{[GM$|$GM]}}
                \texttt{Straf: (5,4) is niet leeg, maar bevat '─'.} \\
\\ 
\texttt{Jullie hebben gezamenlijk 12 van 16 straffen verzameld.} \\
\texttt{Probeer het opnieuw!} \\
            }
        }
    }
    & & \\ \\

    \theutterance \stepcounter{utterance}  
    & & \multicolumn{4}{p{0.6\linewidth}}{
        \cellcolor[rgb]{0.9,0.9,0.9}{
            \makecell[{{p{\linewidth}}}]{
                \texttt{\tiny{[GM$\rangle$P2]}}
                \texttt{Straf: (5,4) is niet leeg, maar bevat '─'.} \\
\\ 
\texttt{Jullie hebben gezamenlijk 12 van 16 straffen verzameld.} \\
\texttt{Probeer het opnieuw!} \\
            }
        }
    }
    & \\ \\

    \theutterance \stepcounter{utterance}  
    & & & \multicolumn{4}{p{0.6\linewidth}}{
        \cellcolor[rgb]{0.9,0.9,0.9}{
            \makecell[{{p{\linewidth}}}]{
                \texttt{\tiny{[GM$\langle$P2]}}
                \texttt{BEWEEG: M, (4, 3)} \\
            }
        }
    }
     \\ \\

    \theutterance \stepcounter{utterance}  
    & & & \multicolumn{2}{p{0.3\linewidth}}{
        \cellcolor[rgb]{0.9,0.9,0.9}{
            \makecell[{{p{\linewidth}}}]{
                \texttt{\tiny{[GM$|$GM]}}
                \texttt{BEWEEG: M, (4, 3)} \\
            }
        }
    }
    & & \\ \\

    \theutterance \stepcounter{utterance}  
    & & & \multicolumn{2}{p{0.3\linewidth}}{
        \cellcolor[rgb]{0.9,0.9,0.9}{
            \makecell[{{p{\linewidth}}}]{
                \texttt{\tiny{[GM$|$GM]}}
                \texttt{Straf: (4,3) is niet leeg, maar bevat 'M'.} \\
\\ 
\texttt{Jullie hebben gezamenlijk 13 van 16 straffen verzameld.} \\
\texttt{Probeer het opnieuw!} \\
            }
        }
    }
    & & \\ \\

    \theutterance \stepcounter{utterance}  
    & & \multicolumn{4}{p{0.6\linewidth}}{
        \cellcolor[rgb]{0.9,0.9,0.9}{
            \makecell[{{p{\linewidth}}}]{
                \texttt{\tiny{[GM$\rangle$P2]}}
                \texttt{Straf: (4,3) is niet leeg, maar bevat 'M'.} \\
\\ 
\texttt{Jullie hebben gezamenlijk 13 van 16 straffen verzameld.} \\
\texttt{Probeer het opnieuw!} \\
            }
        }
    }
    & \\ \\

    \theutterance \stepcounter{utterance}  
    & & & \multicolumn{4}{p{0.6\linewidth}}{
        \cellcolor[rgb]{0.9,0.9,0.9}{
            \makecell[{{p{\linewidth}}}]{
                \texttt{\tiny{[GM$\langle$P2]}}
                \texttt{BEWEEG: O, (5, 4)} \\
            }
        }
    }
     \\ \\

    \theutterance \stepcounter{utterance}  
    & & & \multicolumn{2}{p{0.3\linewidth}}{
        \cellcolor[rgb]{0.9,0.9,0.9}{
            \makecell[{{p{\linewidth}}}]{
                \texttt{\tiny{[GM$|$GM]}}
                \texttt{BEWEEG: O, (5, 4)} \\
            }
        }
    }
    & & \\ \\

    \theutterance \stepcounter{utterance}  
    & & & \multicolumn{2}{p{0.3\linewidth}}{
        \cellcolor[rgb]{0.9,0.9,0.9}{
            \makecell[{{p{\linewidth}}}]{
                \texttt{\tiny{[GM$|$GM]}}
                \texttt{Straf: (5,4) is niet leeg, maar bevat '─'.} \\
\\ 
\texttt{Jullie hebben gezamenlijk 14 van 16 straffen verzameld.} \\
\texttt{Probeer het opnieuw!} \\
            }
        }
    }
    & & \\ \\

    \theutterance \stepcounter{utterance}  
    & & \multicolumn{4}{p{0.6\linewidth}}{
        \cellcolor[rgb]{0.9,0.9,0.9}{
            \makecell[{{p{\linewidth}}}]{
                \texttt{\tiny{[GM$\rangle$P2]}}
                \texttt{Straf: (5,4) is niet leeg, maar bevat '─'.} \\
\\ 
\texttt{Jullie hebben gezamenlijk 14 van 16 straffen verzameld.} \\
\texttt{Probeer het opnieuw!} \\
            }
        }
    }
    & \\ \\

    \theutterance \stepcounter{utterance}  
    & & & \multicolumn{4}{p{0.6\linewidth}}{
        \cellcolor[rgb]{0.9,0.9,0.9}{
            \makecell[{{p{\linewidth}}}]{
                \texttt{\tiny{[GM$\langle$P2]}}
                \texttt{ZEG: Onze doelcellen (4,3), (4,4), (5,4) zijn geblokkeerd door vaste elementen (M, lijnen). Kunnen we de doelen verschuiven naar vrije cellen? Bijv: P(3,3), O(3,4), T(5,3), A(4,3), D(4,4), M(5,3), S(4,5) {-} of andere vrije plekken?} \\
            }
        }
    }
     \\ \\

    \theutterance \stepcounter{utterance}  
    & & & \multicolumn{2}{p{0.3\linewidth}}{
        \cellcolor[rgb]{0.9,0.9,0.9}{
            \makecell[{{p{\linewidth}}}]{
                \texttt{\tiny{[GM$|$GM]}}
                \texttt{ZEG: Onze doelcellen (4,3), (4,4), (5,4) zijn geblokkeerd door vaste elementen (M, lijnen). Kunnen we de doelen verschuiven naar vrije cellen? Bijv: P(3,3), O(3,4), T(5,3), A(4,3), D(4,4), M(5,3), S(4,5) {-} of andere vrije plekken?} \\
            }
        }
    }
    & & \\ \\

    \theutterance \stepcounter{utterance}  
    & & \multicolumn{4}{p{0.6\linewidth}}{
        \cellcolor[rgb]{0.9,0.9,0.9}{
            \makecell[{{p{\linewidth}}}]{
                \texttt{\tiny{[P1$\langle$GM]}}
                \texttt{Verplaatst 'T' naar (5,3) met succes. Uw bijgewerkte raster ziet er als volgt uit:} \\
\texttt{```} \\
\texttt{ 1234567} \\
\texttt{╔═══════╗} \\
\texttt{║◌◌◌◌◌◌◌║ 1} \\
\texttt{║◌◌◌◌◌◌◌║ 2} \\
\texttt{║◌◌PDT┌─╢ 3} \\
\texttt{║◌◌A┌─┤S║ 4} \\
\texttt{║O◌◌│◌├─╢ 5} \\
\texttt{╟───┼─┘◌║ 6} \\
\texttt{║◌M◌│◌◌◌║ 7} \\
\texttt{╚═══╧═══╝} \\
\texttt{```} \\
\texttt{Je speelt momenteel ronde 5 van maximaal 28.} \\
\texttt{Jullie hebben gezamenlijk 14 van 16 straffen verzameld.} \\
\texttt{De andere speler stuurde dit bericht:} \\
\\ 
\\ 
\texttt{"Onze doelcellen (4,3), (4,4), (5,4) zijn geblokkeerd door vaste elementen (M, lijnen). Kunnen we de doelen verschuiven naar vrije cellen? Bijv: P(3,3), O(3,4), T(5,3), A(4,3), D(4,4), M(5,3), S(4,5) {-} of andere vrije plekken?"} \\
\\ 
\\ 
\texttt{Wat is je volgende opdracht?} \\
            }
        }
    }
    & \\ \\

    \theutterance \stepcounter{utterance}  
    & \multicolumn{4}{p{0.6\linewidth}}{
        \cellcolor[rgb]{0.9,0.9,0.9}{
            \makecell[{{p{\linewidth}}}]{
                \texttt{\tiny{[P1$\rangle$GM]}}
                \texttt{ZEG: Onze doel (4,4) en (4,5) zitten ook vast (wand '┌' en lijn '│'). Mijn Vorschlag: P(3,3){[}OK{]}, O(3,4){[}leeg{]}, T(5,3){[}OK{]}, A(3,4){[}bezet door A{]}, D(5,4){[}leeg{]}, M(3,5){[}leeg{]}, S(5,5){[}leeg{]}. Stem je toe met: P(3,3), O(3,4), T(5,3), D(5,4), M(3,5), S(5,5)? A blijft op (3,4).} \\
            }
        }
    }
    & & \\ \\

    \theutterance \stepcounter{utterance}  
    & & & \multicolumn{2}{p{0.3\linewidth}}{
        \cellcolor[rgb]{0.9,0.9,0.9}{
            \makecell[{{p{\linewidth}}}]{
                \texttt{\tiny{[GM$|$GM]}}
                \texttt{ZEG: Onze doel (4,4) en (4,5) zitten ook vast (wand '┌' en lijn '│'). Mijn Vorschlag: P(3,3){[}OK{]}, O(3,4){[}leeg{]}, T(5,3){[}OK{]}, A(3,4){[}bezet door A{]}, D(5,4){[}leeg{]}, M(3,5){[}leeg{]}, S(5,5){[}leeg{]}. Stem je toe met: P(3,3), O(3,4), T(5,3), D(5,4), M(3,5), S(5,5)? A blijft op (3,4).} \\
            }
        }
    }
    & & \\ \\

    \theutterance \stepcounter{utterance}  
    & & \multicolumn{4}{p{0.6\linewidth}}{
        \cellcolor[rgb]{0.9,0.9,0.9}{
            \makecell[{{p{\linewidth}}}]{
                \texttt{\tiny{[GM$\rangle$P2]}}
                \texttt{Je bericht is doorgestuurd naar de andere speler.} \\
\\ 
\\ 
\texttt{Je speelt momenteel ronde 6 van maximaal 28.} \\
\texttt{Jullie hebben gezamenlijk 14 van 16 straffen verzameld.} \\
\texttt{De andere speler stuurde dit bericht:} \\
\\ 
\\ 
\texttt{"Onze doel (4,4) en (4,5) zitten ook vast (wand '┌' en lijn '│'). Mijn Vorschlag: P(3,3){[}OK{]}, O(3,4){[}leeg{]}, T(5,3){[}OK{]}, A(3,4){[}bezet door A{]}, D(5,4){[}leeg{]}, M(3,5){[}leeg{]}, S(5,5){[}leeg{]}. Stem je toe met: P(3,3), O(3,4), T(5,3), D(5,4), M(3,5), S(5,5)? A blijft op (3,4)."} \\
\\ 
\\ 
\texttt{Wat is je volgende opdracht?} \\
            }
        }
    }
    & \\ \\

    \theutterance \stepcounter{utterance}  
    & & & \multicolumn{4}{p{0.6\linewidth}}{
        \cellcolor[rgb]{0.9,0.9,0.9}{
            \makecell[{{p{\linewidth}}}]{
                \texttt{\tiny{[GM$\langle$P2]}}
                \texttt{ZEG: Akkoord met nieuwe doelen: P(3,3){[}OK{]}, O(3,4), T(5,3){[}OK{]}, D(5,4), M(3,5), S(5,5). A blijft op (3,4). Ik ga O, D, M, S verplaatsen.} \\
            }
        }
    }
     \\ \\

    \theutterance \stepcounter{utterance}  
    & & & \multicolumn{2}{p{0.3\linewidth}}{
        \cellcolor[rgb]{0.9,0.9,0.9}{
            \makecell[{{p{\linewidth}}}]{
                \texttt{\tiny{[GM$|$GM]}}
                \texttt{ZEG: Akkoord met nieuwe doelen: P(3,3){[}OK{]}, O(3,4), T(5,3){[}OK{]}, D(5,4), M(3,5), S(5,5). A blijft op (3,4). Ik ga O, D, M, S verplaatsen.} \\
            }
        }
    }
    & & \\ \\

    \theutterance \stepcounter{utterance}  
    & & \multicolumn{4}{p{0.6\linewidth}}{
        \cellcolor[rgb]{0.9,0.9,0.9}{
            \makecell[{{p{\linewidth}}}]{
                \texttt{\tiny{[P1$\langle$GM]}}
                \texttt{Je bericht is doorgestuurd naar de andere speler.} \\
\\ 
\\ 
\texttt{Je speelt momenteel ronde 6 van maximaal 28.} \\
\texttt{Jullie hebben gezamenlijk 14 van 16 straffen verzameld.} \\
\texttt{De andere speler stuurde dit bericht:} \\
\\ 
\\ 
\texttt{"Akkoord met nieuwe doelen: P(3,3){[}OK{]}, O(3,4), T(5,3){[}OK{]}, D(5,4), M(3,5), S(5,5). A blijft op (3,4). Ik ga O, D, M, S verplaatsen."} \\
\\ 
\\ 
\texttt{Wat is je volgende opdracht?} \\
            }
        }
    }
    & \\ \\

    \theutterance \stepcounter{utterance}  
    & \multicolumn{4}{p{0.6\linewidth}}{
        \cellcolor[rgb]{0.9,0.9,0.9}{
            \makecell[{{p{\linewidth}}}]{
                \texttt{\tiny{[P1$\rangle$GM]}}
                \texttt{BEWEEG: O, (3, 4)} \\
            }
        }
    }
    & & \\ \\

    \theutterance \stepcounter{utterance}  
    & & & \multicolumn{2}{p{0.3\linewidth}}{
        \cellcolor[rgb]{0.9,0.9,0.9}{
            \makecell[{{p{\linewidth}}}]{
                \texttt{\tiny{[GM$|$GM]}}
                \texttt{BEWEEG: O, (3, 4)} \\
            }
        }
    }
    & & \\ \\

    \theutterance \stepcounter{utterance}  
    & & & \multicolumn{2}{p{0.3\linewidth}}{
        \cellcolor[rgb]{0.9,0.9,0.9}{
            \makecell[{{p{\linewidth}}}]{
                \texttt{\tiny{[GM$|$GM]}}
                \texttt{Straf: (3,4) is niet leeg, maar bevat 'A'.} \\
\\ 
\texttt{Jullie hebben gezamenlijk 15 van 16 straffen verzameld.} \\
\texttt{Probeer het opnieuw!} \\
            }
        }
    }
    & & \\ \\

    \theutterance \stepcounter{utterance}  
    & & \multicolumn{4}{p{0.6\linewidth}}{
        \cellcolor[rgb]{0.9,0.9,0.9}{
            \makecell[{{p{\linewidth}}}]{
                \texttt{\tiny{[P1$\langle$GM]}}
                \texttt{Straf: (3,4) is niet leeg, maar bevat 'A'.} \\
\\ 
\texttt{Jullie hebben gezamenlijk 15 van 16 straffen verzameld.} \\
\texttt{Probeer het opnieuw!} \\
            }
        }
    }
    & \\ \\

    \theutterance \stepcounter{utterance}  
    & \multicolumn{4}{p{0.6\linewidth}}{
        \cellcolor[rgb]{0.9,0.9,0.9}{
            \makecell[{{p{\linewidth}}}]{
                \texttt{\tiny{[P1$\rangle$GM]}}
                \texttt{BEWEEG: M, (3, 5)} \\
            }
        }
    }
    & & \\ \\

    \theutterance \stepcounter{utterance}  
    & & & \multicolumn{2}{p{0.3\linewidth}}{
        \cellcolor[rgb]{0.9,0.9,0.9}{
            \makecell[{{p{\linewidth}}}]{
                \texttt{\tiny{[GM$|$GM]}}
                \texttt{BEWEEG: M, (3, 5)} \\
            }
        }
    }
    & & \\ \\

    \theutterance \stepcounter{utterance}  
    & & & \multicolumn{2}{p{0.3\linewidth}}{
        \cellcolor[rgb]{0.9,0.9,0.9}{
            \makecell[{{p{\linewidth}}}]{
                \texttt{\tiny{[GM$|$GM]}}
                \texttt{M moved to (3, 5)} \\
            }
        }
    }
    & & \\ \\

    \theutterance \stepcounter{utterance}  
    & & \multicolumn{4}{p{0.6\linewidth}}{
        \cellcolor[rgb]{0.9,0.9,0.9}{
            \makecell[{{p{\linewidth}}}]{
                \texttt{\tiny{[GM$\rangle$P2]}}
                \texttt{Je bericht is doorgestuurd naar de andere speler.} \\
\\ 
\\ 
\texttt{Je speelt momenteel ronde 7 van maximaal 28.} \\
\texttt{Jullie hebben gezamenlijk 15 van 16 straffen verzameld.} \\
\texttt{De andere speler heeft een voorwerp op zijn rooster verplaatst.} \\
\\ 
\\ 
\texttt{Wat is je volgende opdracht?} \\
            }
        }
    }
    & \\ \\

    \theutterance \stepcounter{utterance}  
    & & & \multicolumn{4}{p{0.6\linewidth}}{
        \cellcolor[rgb]{0.9,0.9,0.9}{
            \makecell[{{p{\linewidth}}}]{
                \texttt{\tiny{[GM$\langle$P2]}}
                \texttt{BEWEEG: O, (3, 4)} \\
            }
        }
    }
     \\ \\

    \theutterance \stepcounter{utterance}  
    & & & \multicolumn{2}{p{0.3\linewidth}}{
        \cellcolor[rgb]{0.9,0.9,0.9}{
            \makecell[{{p{\linewidth}}}]{
                \texttt{\tiny{[GM$|$GM]}}
                \texttt{BEWEEG: O, (3, 4)} \\
            }
        }
    }
    & & \\ \\

    \theutterance \stepcounter{utterance}  
    & & & \multicolumn{2}{p{0.3\linewidth}}{
        \cellcolor[rgb]{0.9,0.9,0.9}{
            \makecell[{{p{\linewidth}}}]{
                \texttt{\tiny{[GM$|$GM]}}
                \texttt{Straf: (3,4) is niet leeg, maar bevat 'A'.} \\
\\ 
\texttt{Jullie hebben gezamenlijk 16 van 16 straffen verzameld.} \\
\texttt{Probeer het opnieuw!} \\
            }
        }
    }
    & & \\ \\

    \theutterance \stepcounter{utterance}  
    & & \multicolumn{4}{p{0.6\linewidth}}{
        \cellcolor[rgb]{0.9,0.9,0.9}{
            \makecell[{{p{\linewidth}}}]{
                \texttt{\tiny{[GM$\rangle$P2]}}
                \texttt{Straf: (3,4) is niet leeg, maar bevat 'A'.} \\
\\ 
\texttt{Jullie hebben gezamenlijk 16 van 16 straffen verzameld.} \\
\texttt{Probeer het opnieuw!} \\
            }
        }
    }
    & \\ \\

    \theutterance \stepcounter{utterance}  
    & & & \multicolumn{4}{p{0.6\linewidth}}{
        \cellcolor[rgb]{0.9,0.9,0.9}{
            \makecell[{{p{\linewidth}}}]{
                \texttt{\tiny{[GM$\langle$P2]}}
                \texttt{ZEG: Fout in mijn vorige bericht: A staat al op (3,4). Dus O kan daar niet. Nieuw voorstel: O naar (4,3)? Of (3,5)? Wat is er vrij bij jou? Ik probeer M naar (3,5) en S naar (5,5) en D naar (5,4).} \\
            }
        }
    }
     \\ \\

    \theutterance \stepcounter{utterance}  
    & & & \multicolumn{2}{p{0.3\linewidth}}{
        \cellcolor[rgb]{0.9,0.9,0.9}{
            \makecell[{{p{\linewidth}}}]{
                \texttt{\tiny{[GM$|$GM]}}
                \texttt{ZEG: Fout in mijn vorige bericht: A staat al op (3,4). Dus O kan daar niet. Nieuw voorstel: O naar (4,3)? Of (3,5)? Wat is er vrij bij jou? Ik probeer M naar (3,5) en S naar (5,5) en D naar (5,4).} \\
            }
        }
    }
    & & \\ \\

    \theutterance \stepcounter{utterance}  
    & & \multicolumn{4}{p{0.6\linewidth}}{
        \cellcolor[rgb]{0.9,0.9,0.9}{
            \makecell[{{p{\linewidth}}}]{
                \texttt{\tiny{[P1$\langle$GM]}}
                \texttt{Verplaatst 'M' naar (3,5) met succes. Uw bijgewerkte raster ziet er als volgt uit:} \\
\texttt{```} \\
\texttt{ 1234567} \\
\texttt{╔═══════╗} \\
\texttt{║◌◌◌◌◌◌◌║ 1} \\
\texttt{║◌◌◌◌◌◌◌║ 2} \\
\texttt{║◌◌PDT┌─╢ 3} \\
\texttt{║◌◌A┌─┤S║ 4} \\
\texttt{║O◌M│◌├─╢ 5} \\
\texttt{╟───┼─┘◌║ 6} \\
\texttt{║◌◌◌│◌◌◌║ 7} \\
\texttt{╚═══╧═══╝} \\
\texttt{```} \\
\texttt{Je speelt momenteel ronde 7 van maximaal 28.} \\
\texttt{Jullie hebben gezamenlijk 16 van 16 straffen verzameld.} \\
\texttt{De andere speler stuurde dit bericht:} \\
\\ 
\\ 
\texttt{"Fout in mijn vorige bericht: A staat al op (3,4). Dus O kan daar niet. Nieuw voorstel: O naar (4,3)? Of (3,5)? Wat is er vrij bij jou? Ik probeer M naar (3,5) en S naar (5,5) en D naar (5,4)."} \\
\\ 
\\ 
\texttt{Wat is je volgende opdracht?} \\
            }
        }
    }
    & \\ \\

    \theutterance \stepcounter{utterance}  
    & \multicolumn{4}{p{0.6\linewidth}}{
        \cellcolor[rgb]{0.9,0.9,0.9}{
            \makecell[{{p{\linewidth}}}]{
                \texttt{\tiny{[P1$\rangle$GM]}}
                \texttt{ZEG: (3,5) is nu bezet door M. Vrije cellen in doelgebied: O:(4,3), D:(5,4), S:(5,5). A blijft (3,4), P(3,3), T(5,3), M(3,5). Ga voor O(4,3), D(5,4), S(5,5). Akkoord?} \\
            }
        }
    }
    & & \\ \\

    \theutterance \stepcounter{utterance}  
    & & & \multicolumn{2}{p{0.3\linewidth}}{
        \cellcolor[rgb]{0.9,0.9,0.9}{
            \makecell[{{p{\linewidth}}}]{
                \texttt{\tiny{[GM$|$GM]}}
                \texttt{ZEG: (3,5) is nu bezet door M. Vrije cellen in doelgebied: O:(4,3), D:(5,4), S:(5,5). A blijft (3,4), P(3,3), T(5,3), M(3,5). Ga voor O(4,3), D(5,4), S(5,5). Akkoord?} \\
            }
        }
    }
    & & \\ \\

    \theutterance \stepcounter{utterance}  
    & & \multicolumn{4}{p{0.6\linewidth}}{
        \cellcolor[rgb]{0.9,0.9,0.9}{
            \makecell[{{p{\linewidth}}}]{
                \texttt{\tiny{[GM$\rangle$P2]}}
                \texttt{Je bericht is doorgestuurd naar de andere speler.} \\
\\ 
\\ 
\texttt{Je speelt momenteel ronde 8 van maximaal 28.} \\
\texttt{Jullie hebben gezamenlijk 16 van 16 straffen verzameld.} \\
\texttt{De andere speler stuurde dit bericht:} \\
\\ 
\\ 
\texttt{"(3,5) is nu bezet door M. Vrije cellen in doelgebied: O:(4,3), D:(5,4), S:(5,5). A blijft (3,4), P(3,3), T(5,3), M(3,5). Ga voor O(4,3), D(5,4), S(5,5). Akkoord?"} \\
\\ 
\\ 
\texttt{Wat is je volgende opdracht?} \\
            }
        }
    }
    & \\ \\

    \theutterance \stepcounter{utterance}  
    & & & \multicolumn{4}{p{0.6\linewidth}}{
        \cellcolor[rgb]{0.9,0.9,0.9}{
            \makecell[{{p{\linewidth}}}]{
                \texttt{\tiny{[GM$\langle$P2]}}
                \texttt{BEWEEG: O, (4, 3)} \\
            }
        }
    }
     \\ \\

    \theutterance \stepcounter{utterance}  
    & & & \multicolumn{2}{p{0.3\linewidth}}{
        \cellcolor[rgb]{0.9,0.9,0.9}{
            \makecell[{{p{\linewidth}}}]{
                \texttt{\tiny{[GM$|$GM]}}
                \texttt{BEWEEG: O, (4, 3)} \\
            }
        }
    }
    & & \\ \\

    \theutterance \stepcounter{utterance}  
    & & & \multicolumn{2}{p{0.3\linewidth}}{
        \cellcolor[rgb]{0.9,0.9,0.9}{
            \makecell[{{p{\linewidth}}}]{
                \texttt{\tiny{[GM$|$GM]}}
                \texttt{Straf: (4,3) is niet leeg, maar bevat 'M'.} \\
\\ 
\texttt{Jullie hebben gezamenlijk 17 van 16 straffen verzameld.} \\
\texttt{Probeer het opnieuw!} \\
            }
        }
    }
    & & \\ \\

    \theutterance \stepcounter{utterance}  
    & & \multicolumn{4}{p{0.6\linewidth}}{
        \cellcolor[rgb]{0.9,0.9,0.9}{
            \makecell[{{p{\linewidth}}}]{
                \texttt{\tiny{[GM$\rangle$P2]}}
                \texttt{Straf: (4,3) is niet leeg, maar bevat 'M'.} \\
\\ 
\texttt{Jullie hebben gezamenlijk 17 van 16 straffen verzameld.} \\
\texttt{Probeer het opnieuw!} \\
            }
        }
    }
    & \\ \\

    \theutterance \stepcounter{utterance}  
    & & & \multicolumn{4}{p{0.6\linewidth}}{
        \cellcolor[rgb]{0.9,0.9,0.9}{
            \makecell[{{p{\linewidth}}}]{
                \texttt{\tiny{[GM$\langle$P2]}}
                \texttt{BEWEEG: M, (3, 5)} \\
            }
        }
    }
     \\ \\

    \theutterance \stepcounter{utterance}  
    & & & \multicolumn{2}{p{0.3\linewidth}}{
        \cellcolor[rgb]{0.9,0.9,0.9}{
            \makecell[{{p{\linewidth}}}]{
                \texttt{\tiny{[GM$|$GM]}}
                \texttt{BEWEEG: M, (3, 5)} \\
            }
        }
    }
    & & \\ \\

    \theutterance \stepcounter{utterance}  
    & & & \multicolumn{2}{p{0.3\linewidth}}{
        \cellcolor[rgb]{0.9,0.9,0.9}{
            \makecell[{{p{\linewidth}}}]{
                \texttt{\tiny{[GM$|$GM]}}
                \texttt{M moved to (3, 5)} \\
            }
        }
    }
    & & \\ \\

    \theutterance \stepcounter{utterance}  
    & & \multicolumn{4}{p{0.6\linewidth}}{
        \cellcolor[rgb]{0.9,0.9,0.9}{
            \makecell[{{p{\linewidth}}}]{
                \texttt{\tiny{[P1$\langle$GM]}}
                \texttt{Je bericht is doorgestuurd naar de andere speler.} \\
\\ 
\\ 
\texttt{Je speelt momenteel ronde 8 van maximaal 28.} \\
\texttt{Jullie hebben gezamenlijk 17 van 16 straffen verzameld.} \\
\texttt{De andere speler heeft een voorwerp op zijn rooster verplaatst.} \\
\\ 
\\ 
\texttt{Wat is je volgende opdracht?} \\
            }
        }
    }
    & \\ \\

    \theutterance \stepcounter{utterance}  
    & \multicolumn{4}{p{0.6\linewidth}}{
        \cellcolor[rgb]{0.9,0.9,0.9}{
            \makecell[{{p{\linewidth}}}]{
                \texttt{\tiny{[P1$\rangle$GM]}}
                \texttt{ZEG: Onze straffen zijn op (17/16), we mogen geen fouten meer maken. Controleer of (4,3), (5,4), (5,5) leeg zijn bij jou voordat je zet. Ik wacht op je bevestiging of succesvoeling voordat ik mijn O, D, S verplaats.} \\
            }
        }
    }
    & & \\ \\

    \theutterance \stepcounter{utterance}  
    & & & \multicolumn{2}{p{0.3\linewidth}}{
        \cellcolor[rgb]{0.9,0.9,0.9}{
            \makecell[{{p{\linewidth}}}]{
                \texttt{\tiny{[GM$|$GM]}}
                \texttt{ZEG: Onze straffen zijn op (17/16), we mogen geen fouten meer maken. Controleer of (4,3), (5,4), (5,5) leeg zijn bij jou voordat je zet. Ik wacht op je bevestiging of succesvoeling voordat ik mijn O, D, S verplaats.} \\
            }
        }
    }
    & & \\ \\

    \theutterance \stepcounter{utterance}  
    & & \multicolumn{4}{p{0.6\linewidth}}{
        \cellcolor[rgb]{0.9,0.9,0.9}{
            \makecell[{{p{\linewidth}}}]{
                \texttt{\tiny{[GM$\rangle$P2]}}
                \texttt{Verplaatst 'M' naar (3,5) met succes. Uw bijgewerkte raster ziet er als volgt uit:} \\
\texttt{```} \\
\texttt{ 1234567} \\
\texttt{╔═══════╗} \\
\texttt{║◌◌◌◌◌S◌║ 1} \\
\texttt{║◌◌◌◌O◌◌║ 2} \\
\texttt{║◌◌P◌T┌─╢ 3} \\
\texttt{║◌◌A┌─┤◌║ 4} \\
\texttt{║◌◌M│◌├─╢ 5} \\
\texttt{╟───┼─┘◌║ 6} \\
\texttt{║D◌◌│◌◌◌║ 7} \\
\texttt{╚═══╧═══╝} \\
\texttt{```} \\
\texttt{Je speelt momenteel ronde 9 van maximaal 28.} \\
\texttt{Jullie hebben gezamenlijk 17 van 16 straffen verzameld.} \\
\texttt{De andere speler stuurde dit bericht:} \\
\\ 
\\ 
\texttt{"Onze straffen zijn op (17/16), we mogen geen fouten meer maken. Controleer of (4,3), (5,4), (5,5) leeg zijn bij jou voordat je zet. Ik wacht op je bevestiging of succesvoeling voordat ik mijn O, D, S verplaats."} \\
\\ 
\\ 
\texttt{Wat is je volgende opdracht?} \\
            }
        }
    }
    & \\ \\

    \theutterance \stepcounter{utterance}  
    & & & \multicolumn{4}{p{0.6\linewidth}}{
        \cellcolor[rgb]{0.9,0.9,0.9}{
            \makecell[{{p{\linewidth}}}]{
                \texttt{\tiny{[GM$\langle$P2]}}
                \texttt{ZEG: Mijn vrije cellen in doelgebied: (4,3) bevat 'M', (5,4) bevat '─' (lijn), (5,5) leeg. Alleen (5,5) is vrij. Kunnen we O en D elders plaatsen? Bijv. O(3,5) is nu bezet door M. O(2,4)? D(4,5)? S(5,5) gaat lukken. Wat zijn jullie vrije cellen?} \\
            }
        }
    }
     \\ \\

    \theutterance \stepcounter{utterance}  
    & & & \multicolumn{2}{p{0.3\linewidth}}{
        \cellcolor[rgb]{0.9,0.9,0.9}{
            \makecell[{{p{\linewidth}}}]{
                \texttt{\tiny{[GM$|$GM]}}
                \texttt{ZEG: Mijn vrije cellen in doelgebied: (4,3) bevat 'M', (5,4) bevat '─' (lijn), (5,5) leeg. Alleen (5,5) is vrij. Kunnen we O en D elders plaatsen? Bijv. O(3,5) is nu bezet door M. O(2,4)? D(4,5)? S(5,5) gaat lukken. Wat zijn jullie vrije cellen?} \\
            }
        }
    }
    & & \\ \\

    \theutterance \stepcounter{utterance}  
    & & \multicolumn{4}{p{0.6\linewidth}}{
        \cellcolor[rgb]{0.9,0.9,0.9}{
            \makecell[{{p{\linewidth}}}]{
                \texttt{\tiny{[P1$\langle$GM]}}
                \texttt{Je bericht is doorgestuurd naar de andere speler.} \\
\\ 
\\ 
\texttt{Je speelt momenteel ronde 9 van maximaal 28.} \\
\texttt{Jullie hebben gezamenlijk 17 van 16 straffen verzameld.} \\
\texttt{De andere speler stuurde dit bericht:} \\
\\ 
\\ 
\texttt{"Mijn vrije cellen in doelgebied: (4,3) bevat 'M', (5,4) bevat '─' (lijn), (5,5) leeg. Alleen (5,5) is vrij. Kunnen we O en D elders plaatsen? Bijv. O(3,5) is nu bezet door M. O(2,4)? D(4,5)? S(5,5) gaat lukken. Wat zijn jullie vrije cellen?"} \\
\\ 
\\ 
\texttt{Wat is je volgende opdracht?} \\
            }
        }
    }
    & \\ \\

    \theutterance \stepcounter{utterance}  
    & \multicolumn{4}{p{0.6\linewidth}}{
        \cellcolor[rgb]{0.9,0.9,0.9}{
            \makecell[{{p{\linewidth}}}]{
                \texttt{\tiny{[P1$\rangle$GM]}}
                \texttt{ZEG: Mijn vrije cellen in nabijheid: (4,3) leeg, (4,5) muur, (5,4) leeg, (5,5) leeg, (2,4) leeg, (3,6) leeg. Voorstel nieuwe doelen (alleen lege cellen): O(4,3), D(5,4), S(5,5). P(3,3), T(5,3), A(3,4), M(3,5) zitten al goed. Stem je toe?} \\
            }
        }
    }
    & & \\ \\

    \theutterance \stepcounter{utterance}  
    & & & \multicolumn{2}{p{0.3\linewidth}}{
        \cellcolor[rgb]{0.9,0.9,0.9}{
            \makecell[{{p{\linewidth}}}]{
                \texttt{\tiny{[GM$|$GM]}}
                \texttt{ZEG: Mijn vrije cellen in nabijheid: (4,3) leeg, (4,5) muur, (5,4) leeg, (5,5) leeg, (2,4) leeg, (3,6) leeg. Voorstel nieuwe doelen (alleen lege cellen): O(4,3), D(5,4), S(5,5). P(3,3), T(5,3), A(3,4), M(3,5) zitten al goed. Stem je toe?} \\
            }
        }
    }
    & & \\ \\

    \theutterance \stepcounter{utterance}  
    & & \multicolumn{4}{p{0.6\linewidth}}{
        \cellcolor[rgb]{0.9,0.9,0.9}{
            \makecell[{{p{\linewidth}}}]{
                \texttt{\tiny{[GM$\rangle$P2]}}
                \texttt{Je bericht is doorgestuurd naar de andere speler.} \\
\\ 
\\ 
\texttt{Je speelt momenteel ronde 10 van maximaal 28.} \\
\texttt{Jullie hebben gezamenlijk 17 van 16 straffen verzameld.} \\
\texttt{De andere speler stuurde dit bericht:} \\
\\ 
\\ 
\texttt{"Mijn vrije cellen in nabijheid: (4,3) leeg, (4,5) muur, (5,4) leeg, (5,5) leeg, (2,4) leeg, (3,6) leeg. Voorstel nieuwe doelen (alleen lege cellen): O(4,3), D(5,4), S(5,5). P(3,3), T(5,3), A(3,4), M(3,5) zitten al goed. Stem je toe?"} \\
\\ 
\\ 
\texttt{Wat is je volgende opdracht?} \\
            }
        }
    }
    & \\ \\

    \theutterance \stepcounter{utterance}  
    & & & \multicolumn{4}{p{0.6\linewidth}}{
        \cellcolor[rgb]{0.9,0.9,0.9}{
            \makecell[{{p{\linewidth}}}]{
                \texttt{\tiny{[GM$\langle$P2]}}
                \texttt{ZEG: Mijn (4,3) zit vol met 'M' (staat op 4,3), (5,4) is lijn '─', (5,5) leeg. Alleen S(5,5) matcht. O(4,3) en D(5,4) werken bij mij niet. Kunnen we O en D op (3,6) en (2,4) zetten? Of (4,5) is muur. M(3,5) is bezet. Ik stem toe voor S(5,5). Voor O en D: O(3,6), D(2,4)? Akkoord?} \\
            }
        }
    }
     \\ \\

    \theutterance \stepcounter{utterance}  
    & & & \multicolumn{2}{p{0.3\linewidth}}{
        \cellcolor[rgb]{0.9,0.9,0.9}{
            \makecell[{{p{\linewidth}}}]{
                \texttt{\tiny{[GM$|$GM]}}
                \texttt{ZEG: Mijn (4,3) zit vol met 'M' (staat op 4,3), (5,4) is lijn '─', (5,5) leeg. Alleen S(5,5) matcht. O(4,3) en D(5,4) werken bij mij niet. Kunnen we O en D op (3,6) en (2,4) zetten? Of (4,5) is muur. M(3,5) is bezet. Ik stem toe voor S(5,5). Voor O en D: O(3,6), D(2,4)? Akkoord?} \\
            }
        }
    }
    & & \\ \\

    \theutterance \stepcounter{utterance}  
    & & \multicolumn{4}{p{0.6\linewidth}}{
        \cellcolor[rgb]{0.9,0.9,0.9}{
            \makecell[{{p{\linewidth}}}]{
                \texttt{\tiny{[P1$\langle$GM]}}
                \texttt{Je bericht is doorgestuurd naar de andere speler.} \\
\\ 
\\ 
\texttt{Je speelt momenteel ronde 10 van maximaal 28.} \\
\texttt{Jullie hebben gezamenlijk 17 van 16 straffen verzameld.} \\
\texttt{De andere speler stuurde dit bericht:} \\
\\ 
\\ 
\texttt{"Mijn (4,3) zit vol met 'M' (staat op 4,3), (5,4) is lijn '─', (5,5) leeg. Alleen S(5,5) matcht. O(4,3) en D(5,4) werken bij mij niet. Kunnen we O en D op (3,6) en (2,4) zetten? Of (4,5) is muur. M(3,5) is bezet. Ik stem toe voor S(5,5). Voor O en D: O(3,6), D(2,4)? Akkoord?"} \\
\\ 
\\ 
\texttt{Wat is je volgende opdracht?} \\
            }
        }
    }
    & \\ \\

    \theutterance \stepcounter{utterance}  
    & \multicolumn{4}{p{0.6\linewidth}}{
        \cellcolor[rgb]{0.9,0.9,0.9}{
            \makecell[{{p{\linewidth}}}]{
                \texttt{\tiny{[P1$\rangle$GM]}}
                \texttt{ZEG: Mijn (3,6) en (2,4) zijn beide leeg. Akkoord met O(3,6), D(2,4), S(5,5). Overzicht finale doelen: P(3,3), O(3,6), T(5,3), A(3,4), D(2,4), M(3,5), S(5,5). Ik ga nu O, D, S verplaatsen.} \\
            }
        }
    }
    & & \\ \\

    \theutterance \stepcounter{utterance}  
    & & & \multicolumn{2}{p{0.3\linewidth}}{
        \cellcolor[rgb]{0.9,0.9,0.9}{
            \makecell[{{p{\linewidth}}}]{
                \texttt{\tiny{[GM$|$GM]}}
                \texttt{ZEG: Mijn (3,6) en (2,4) zijn beide leeg. Akkoord met O(3,6), D(2,4), S(5,5). Overzicht finale doelen: P(3,3), O(3,6), T(5,3), A(3,4), D(2,4), M(3,5), S(5,5). Ik ga nu O, D, S verplaatsen.} \\
            }
        }
    }
    & & \\ \\

    \theutterance \stepcounter{utterance}  
    & & \multicolumn{4}{p{0.6\linewidth}}{
        \cellcolor[rgb]{0.9,0.9,0.9}{
            \makecell[{{p{\linewidth}}}]{
                \texttt{\tiny{[GM$\rangle$P2]}}
                \texttt{Je bericht is doorgestuurd naar de andere speler.} \\
\\ 
\\ 
\texttt{Je speelt momenteel ronde 11 van maximaal 28.} \\
\texttt{Jullie hebben gezamenlijk 17 van 16 straffen verzameld.} \\
\texttt{De andere speler stuurde dit bericht:} \\
\\ 
\\ 
\texttt{"Mijn (3,6) en (2,4) zijn beide leeg. Akkoord met O(3,6), D(2,4), S(5,5). Overzicht finale doelen: P(3,3), O(3,6), T(5,3), A(3,4), D(2,4), M(3,5), S(5,5). Ik ga nu O, D, S verplaatsen."} \\
\\ 
\\ 
\texttt{Wat is je volgende opdracht?} \\
            }
        }
    }
    & \\ \\

    \theutterance \stepcounter{utterance}  
    & & & \multicolumn{4}{p{0.6\linewidth}}{
        \cellcolor[rgb]{0.9,0.9,0.9}{
            \makecell[{{p{\linewidth}}}]{
                \texttt{\tiny{[GM$\langle$P2]}}
                \texttt{BEWEEG: S, (5, 5)} \\
            }
        }
    }
     \\ \\

    \theutterance \stepcounter{utterance}  
    & & & \multicolumn{2}{p{0.3\linewidth}}{
        \cellcolor[rgb]{0.9,0.9,0.9}{
            \makecell[{{p{\linewidth}}}]{
                \texttt{\tiny{[GM$|$GM]}}
                \texttt{BEWEEG: S, (5, 5)} \\
            }
        }
    }
    & & \\ \\

    \theutterance \stepcounter{utterance}  
    & & & \multicolumn{2}{p{0.3\linewidth}}{
        \cellcolor[rgb]{0.9,0.9,0.9}{
            \makecell[{{p{\linewidth}}}]{
                \texttt{\tiny{[GM$|$GM]}}
                \texttt{S moved to (5, 5)} \\
            }
        }
    }
    & & \\ \\

    \theutterance \stepcounter{utterance}  
    & & \multicolumn{4}{p{0.6\linewidth}}{
        \cellcolor[rgb]{0.9,0.9,0.9}{
            \makecell[{{p{\linewidth}}}]{
                \texttt{\tiny{[P1$\langle$GM]}}
                \texttt{Je bericht is doorgestuurd naar de andere speler.} \\
\\ 
\\ 
\texttt{Je speelt momenteel ronde 11 van maximaal 28.} \\
\texttt{Jullie hebben gezamenlijk 17 van 16 straffen verzameld.} \\
\texttt{De andere speler heeft een voorwerp op zijn rooster verplaatst.} \\
\\ 
\\ 
\texttt{Wat is je volgende opdracht?} \\
            }
        }
    }
    & \\ \\

    \theutterance \stepcounter{utterance}  
    & \multicolumn{4}{p{0.6\linewidth}}{
        \cellcolor[rgb]{0.9,0.9,0.9}{
            \makecell[{{p{\linewidth}}}]{
                \texttt{\tiny{[P1$\rangle$GM]}}
                \texttt{BEWEEG: O, (3, 6)} \\
            }
        }
    }
    & & \\ \\

    \theutterance \stepcounter{utterance}  
    & & & \multicolumn{2}{p{0.3\linewidth}}{
        \cellcolor[rgb]{0.9,0.9,0.9}{
            \makecell[{{p{\linewidth}}}]{
                \texttt{\tiny{[GM$|$GM]}}
                \texttt{BEWEEG: O, (3, 6)} \\
            }
        }
    }
    & & \\ \\

    \theutterance \stepcounter{utterance}  
    & & & \multicolumn{2}{p{0.3\linewidth}}{
        \cellcolor[rgb]{0.9,0.9,0.9}{
            \makecell[{{p{\linewidth}}}]{
                \texttt{\tiny{[GM$|$GM]}}
                \texttt{Straf: (3,6) is niet leeg, maar bevat '─'.} \\
\\ 
\texttt{Jullie hebben gezamenlijk 18 van 16 straffen verzameld.} \\
\texttt{Probeer het opnieuw!} \\
            }
        }
    }
    & & \\ \\

    \theutterance \stepcounter{utterance}  
    & & \multicolumn{4}{p{0.6\linewidth}}{
        \cellcolor[rgb]{0.9,0.9,0.9}{
            \makecell[{{p{\linewidth}}}]{
                \texttt{\tiny{[P1$\langle$GM]}}
                \texttt{Straf: (3,6) is niet leeg, maar bevat '─'.} \\
\\ 
\texttt{Jullie hebben gezamenlijk 18 van 16 straffen verzameld.} \\
\texttt{Probeer het opnieuw!} \\
            }
        }
    }
    & \\ \\

    \theutterance \stepcounter{utterance}  
    & \multicolumn{4}{p{0.6\linewidth}}{
        \cellcolor[rgb]{0.9,0.9,0.9}{
            \makecell[{{p{\linewidth}}}]{
                \texttt{\tiny{[P1$\rangle$GM]}}
                \texttt{ZEG: Fout, (3,6) is muur. Mijn leegtes: (2,4), (4,3), (5,4), (5,5), (1,3), (1,4), (1,5), (2,5), (2,6). O naar (2,4)? D naar (4,3)? S naar (5,5)? Akkoord?} \\
            }
        }
    }
    & & \\ \\

    \theutterance \stepcounter{utterance}  
    & & & \multicolumn{2}{p{0.3\linewidth}}{
        \cellcolor[rgb]{0.9,0.9,0.9}{
            \makecell[{{p{\linewidth}}}]{
                \texttt{\tiny{[GM$|$GM]}}
                \texttt{ZEG: Fout, (3,6) is muur. Mijn leegtes: (2,4), (4,3), (5,4), (5,5), (1,3), (1,4), (1,5), (2,5), (2,6). O naar (2,4)? D naar (4,3)? S naar (5,5)? Akkoord?} \\
            }
        }
    }
    & & \\ \\

    \theutterance \stepcounter{utterance}  
    & & \multicolumn{4}{p{0.6\linewidth}}{
        \cellcolor[rgb]{0.9,0.9,0.9}{
            \makecell[{{p{\linewidth}}}]{
                \texttt{\tiny{[GM$\rangle$P2]}}
                \texttt{Verplaatst 'S' naar (5,5) met succes. Uw bijgewerkte raster ziet er als volgt uit:} \\
\texttt{```} \\
\texttt{ 1234567} \\
\texttt{╔═══════╗} \\
\texttt{║◌◌◌◌◌◌◌║ 1} \\
\texttt{║◌◌◌◌O◌◌║ 2} \\
\texttt{║◌◌P◌T┌─╢ 3} \\
\texttt{║◌◌A┌─┤◌║ 4} \\
\texttt{║◌◌M│S├─╢ 5} \\
\texttt{╟───┼─┘◌║ 6} \\
\texttt{║D◌◌│◌◌◌║ 7} \\
\texttt{╚═══╧═══╝} \\
\texttt{```} \\
\texttt{Je speelt momenteel ronde 12 van maximaal 28.} \\
\texttt{Jullie hebben gezamenlijk 18 van 16 straffen verzameld.} \\
\texttt{De andere speler stuurde dit bericht:} \\
\\ 
\\ 
\texttt{"Fout, (3,6) is muur. Mijn leegtes: (2,4), (4,3), (5,4), (5,5), (1,3), (1,4), (1,5), (2,5), (2,6). O naar (2,4)? D naar (4,3)? S naar (5,5)? Akkoord?"} \\
\\ 
\\ 
\texttt{Wat is je volgende opdracht?} \\
            }
        }
    }
    & \\ \\

    \theutterance \stepcounter{utterance}  
    & & & \multicolumn{4}{p{0.6\linewidth}}{
        \cellcolor[rgb]{0.9,0.9,0.9}{
            \makecell[{{p{\linewidth}}}]{
                \texttt{\tiny{[GM$\langle$P2]}}
                \texttt{ZEG: Mijn (4,3) zit vol met 'M', (5,4) is lijn. Alleen (2,4) en (5,5) matchen. S(5,5) done. O(2,4) en D(4,3) werkt niet voor mij. D(2,4) kan niet, O(2,4) wel. Kunnen we O(2,4) en D(1,3) of D(1,4) doen? Mijn (1,3), (1,4), (1,5), (2,5), (2,6) zijn leeg. Akkoord met O(2,4), D(1,3)?} \\
            }
        }
    }
     \\ \\

    \theutterance \stepcounter{utterance}  
    & & & \multicolumn{2}{p{0.3\linewidth}}{
        \cellcolor[rgb]{0.9,0.9,0.9}{
            \makecell[{{p{\linewidth}}}]{
                \texttt{\tiny{[GM$|$GM]}}
                \texttt{ZEG: Mijn (4,3) zit vol met 'M', (5,4) is lijn. Alleen (2,4) en (5,5) matchen. S(5,5) done. O(2,4) en D(4,3) werkt niet voor mij. D(2,4) kan niet, O(2,4) wel. Kunnen we O(2,4) en D(1,3) of D(1,4) doen? Mijn (1,3), (1,4), (1,5), (2,5), (2,6) zijn leeg. Akkoord met O(2,4), D(1,3)?} \\
            }
        }
    }
    & & \\ \\

    \theutterance \stepcounter{utterance}  
    & & \multicolumn{4}{p{0.6\linewidth}}{
        \cellcolor[rgb]{0.9,0.9,0.9}{
            \makecell[{{p{\linewidth}}}]{
                \texttt{\tiny{[P1$\langle$GM]}}
                \texttt{Je bericht is doorgestuurd naar de andere speler.} \\
\\ 
\\ 
\texttt{Je speelt momenteel ronde 12 van maximaal 28.} \\
\texttt{Jullie hebben gezamenlijk 18 van 16 straffen verzameld.} \\
\texttt{De andere speler stuurde dit bericht:} \\
\\ 
\\ 
\texttt{"Mijn (4,3) zit vol met 'M', (5,4) is lijn. Alleen (2,4) en (5,5) matchen. S(5,5) done. O(2,4) en D(4,3) werkt niet voor mij. D(2,4) kan niet, O(2,4) wel. Kunnen we O(2,4) en D(1,3) of D(1,4) doen? Mijn (1,3), (1,4), (1,5), (2,5), (2,6) zijn leeg. Akkoord met O(2,4), D(1,3)?"} \\
\\ 
\\ 
\texttt{Wat is je volgende opdracht?} \\
            }
        }
    }
    & \\ \\

    \theutterance \stepcounter{utterance}  
    & \multicolumn{4}{p{0.6\linewidth}}{
        \cellcolor[rgb]{0.9,0.9,0.9}{
            \makecell[{{p{\linewidth}}}]{
                \texttt{\tiny{[P1$\rangle$GM]}}
                \texttt{ZEG: Mijn (1,3) en (1,4) zijn leeg. Akkoord met O(2,4), D(1,3), S(5,5). Finale doelen: P(3,3), O(2,4), T(5,3), A(3,4), D(1,3), M(3,5), S(5,5). Ik zet nu O, D, S.} \\
            }
        }
    }
    & & \\ \\

    \theutterance \stepcounter{utterance}  
    & & & \multicolumn{2}{p{0.3\linewidth}}{
        \cellcolor[rgb]{0.9,0.9,0.9}{
            \makecell[{{p{\linewidth}}}]{
                \texttt{\tiny{[GM$|$GM]}}
                \texttt{ZEG: Mijn (1,3) en (1,4) zijn leeg. Akkoord met O(2,4), D(1,3), S(5,5). Finale doelen: P(3,3), O(2,4), T(5,3), A(3,4), D(1,3), M(3,5), S(5,5). Ik zet nu O, D, S.} \\
            }
        }
    }
    & & \\ \\

    \theutterance \stepcounter{utterance}  
    & & \multicolumn{4}{p{0.6\linewidth}}{
        \cellcolor[rgb]{0.9,0.9,0.9}{
            \makecell[{{p{\linewidth}}}]{
                \texttt{\tiny{[GM$\rangle$P2]}}
                \texttt{Je bericht is doorgestuurd naar de andere speler.} \\
\\ 
\\ 
\texttt{Je speelt momenteel ronde 13 van maximaal 28.} \\
\texttt{Jullie hebben gezamenlijk 18 van 16 straffen verzameld.} \\
\texttt{De andere speler stuurde dit bericht:} \\
\\ 
\\ 
\texttt{"Mijn (1,3) en (1,4) zijn leeg. Akkoord met O(2,4), D(1,3), S(5,5). Finale doelen: P(3,3), O(2,4), T(5,3), A(3,4), D(1,3), M(3,5), S(5,5). Ik zet nu O, D, S."} \\
\\ 
\\ 
\texttt{Wat is je volgende opdracht?} \\
            }
        }
    }
    & \\ \\

    \theutterance \stepcounter{utterance}  
    & & & \multicolumn{4}{p{0.6\linewidth}}{
        \cellcolor[rgb]{0.9,0.9,0.9}{
            \makecell[{{p{\linewidth}}}]{
                \texttt{\tiny{[GM$\langle$P2]}}
                \texttt{BEWEEG: O, (2, 4)} \\
            }
        }
    }
     \\ \\

    \theutterance \stepcounter{utterance}  
    & & & \multicolumn{2}{p{0.3\linewidth}}{
        \cellcolor[rgb]{0.9,0.9,0.9}{
            \makecell[{{p{\linewidth}}}]{
                \texttt{\tiny{[GM$|$GM]}}
                \texttt{BEWEEG: O, (2, 4)} \\
            }
        }
    }
    & & \\ \\

    \theutterance \stepcounter{utterance}  
    & & & \multicolumn{2}{p{0.3\linewidth}}{
        \cellcolor[rgb]{0.9,0.9,0.9}{
            \makecell[{{p{\linewidth}}}]{
                \texttt{\tiny{[GM$|$GM]}}
                \texttt{O moved to (2, 4)} \\
            }
        }
    }
    & & \\ \\

    \theutterance \stepcounter{utterance}  
    & & \multicolumn{4}{p{0.6\linewidth}}{
        \cellcolor[rgb]{0.9,0.9,0.9}{
            \makecell[{{p{\linewidth}}}]{
                \texttt{\tiny{[P1$\langle$GM]}}
                \texttt{Je bericht is doorgestuurd naar de andere speler.} \\
\\ 
\\ 
\texttt{Je speelt momenteel ronde 13 van maximaal 28.} \\
\texttt{Jullie hebben gezamenlijk 18 van 16 straffen verzameld.} \\
\texttt{De andere speler heeft een voorwerp op zijn rooster verplaatst.} \\
\\ 
\\ 
\texttt{Wat is je volgende opdracht?} \\
            }
        }
    }
    & \\ \\

    \theutterance \stepcounter{utterance}  
    & \multicolumn{4}{p{0.6\linewidth}}{
        \cellcolor[rgb]{0.9,0.9,0.9}{
            \makecell[{{p{\linewidth}}}]{
                \texttt{\tiny{[P1$\rangle$GM]}}
                \texttt{BEWEEG: O, (2, 4)} \\
            }
        }
    }
    & & \\ \\

    \theutterance \stepcounter{utterance}  
    & & & \multicolumn{2}{p{0.3\linewidth}}{
        \cellcolor[rgb]{0.9,0.9,0.9}{
            \makecell[{{p{\linewidth}}}]{
                \texttt{\tiny{[GM$|$GM]}}
                \texttt{BEWEEG: O, (2, 4)} \\
            }
        }
    }
    & & \\ \\

    \theutterance \stepcounter{utterance}  
    & & & \multicolumn{2}{p{0.3\linewidth}}{
        \cellcolor[rgb]{0.9,0.9,0.9}{
            \makecell[{{p{\linewidth}}}]{
                \texttt{\tiny{[GM$|$GM]}}
                \texttt{O moved to (2, 4)} \\
            }
        }
    }
    & & \\ \\

    \theutterance \stepcounter{utterance}  
    & & \multicolumn{4}{p{0.6\linewidth}}{
        \cellcolor[rgb]{0.9,0.9,0.9}{
            \makecell[{{p{\linewidth}}}]{
                \texttt{\tiny{[GM$\rangle$P2]}}
                \texttt{Verplaatst 'O' naar (2,4) met succes. Uw bijgewerkte raster ziet er als volgt uit:} \\
\texttt{```} \\
\texttt{ 1234567} \\
\texttt{╔═══════╗} \\
\texttt{║◌◌◌◌◌◌◌║ 1} \\
\texttt{║◌◌◌◌◌◌◌║ 2} \\
\texttt{║◌◌P◌T┌─╢ 3} \\
\texttt{║◌OA┌─┤◌║ 4} \\
\texttt{║◌◌M│S├─╢ 5} \\
\texttt{╟───┼─┘◌║ 6} \\
\texttt{║D◌◌│◌◌◌║ 7} \\
\texttt{╚═══╧═══╝} \\
\texttt{```} \\
\texttt{Je speelt momenteel ronde 14 van maximaal 28.} \\
\texttt{Jullie hebben gezamenlijk 18 van 16 straffen verzameld.} \\
\texttt{De andere speler heeft een voorwerp op zijn rooster verplaatst.} \\
\\ 
\\ 
\texttt{Wat is je volgende opdracht?} \\
            }
        }
    }
    & \\ \\

    \theutterance \stepcounter{utterance}  
    & & & \multicolumn{4}{p{0.6\linewidth}}{
        \cellcolor[rgb]{0.9,0.9,0.9}{
            \makecell[{{p{\linewidth}}}]{
                \texttt{\tiny{[GM$\langle$P2]}}
                \texttt{BEWEEG: D, (1, 3)} \\
            }
        }
    }
     \\ \\

    \theutterance \stepcounter{utterance}  
    & & & \multicolumn{2}{p{0.3\linewidth}}{
        \cellcolor[rgb]{0.9,0.9,0.9}{
            \makecell[{{p{\linewidth}}}]{
                \texttt{\tiny{[GM$|$GM]}}
                \texttt{BEWEEG: D, (1, 3)} \\
            }
        }
    }
    & & \\ \\

    \theutterance \stepcounter{utterance}  
    & & & \multicolumn{2}{p{0.3\linewidth}}{
        \cellcolor[rgb]{0.9,0.9,0.9}{
            \makecell[{{p{\linewidth}}}]{
                \texttt{\tiny{[GM$|$GM]}}
                \texttt{D moved to (1, 3)} \\
            }
        }
    }
    & & \\ \\

    \theutterance \stepcounter{utterance}  
    & & \multicolumn{4}{p{0.6\linewidth}}{
        \cellcolor[rgb]{0.9,0.9,0.9}{
            \makecell[{{p{\linewidth}}}]{
                \texttt{\tiny{[P1$\langle$GM]}}
                \texttt{Verplaatst 'O' naar (2,4) met succes. Uw bijgewerkte raster ziet er als volgt uit:} \\
\texttt{```} \\
\texttt{ 1234567} \\
\texttt{╔═══════╗} \\
\texttt{║◌◌◌◌◌◌◌║ 1} \\
\texttt{║◌◌◌◌◌◌◌║ 2} \\
\texttt{║◌◌PDT┌─╢ 3} \\
\texttt{║◌OA┌─┤S║ 4} \\
\texttt{║◌◌M│◌├─╢ 5} \\
\texttt{╟───┼─┘◌║ 6} \\
\texttt{║◌◌◌│◌◌◌║ 7} \\
\texttt{╚═══╧═══╝} \\
\texttt{```} \\
\texttt{Je speelt momenteel ronde 14 van maximaal 28.} \\
\texttt{Jullie hebben gezamenlijk 18 van 16 straffen verzameld.} \\
\texttt{De andere speler heeft een voorwerp op zijn rooster verplaatst.} \\
\\ 
\\ 
\texttt{Wat is je volgende opdracht?} \\
            }
        }
    }
    & \\ \\

    \theutterance \stepcounter{utterance}  
    & \multicolumn{4}{p{0.6\linewidth}}{
        \cellcolor[rgb]{0.9,0.9,0.9}{
            \makecell[{{p{\linewidth}}}]{
                \texttt{\tiny{[P1$\rangle$GM]}}
                \texttt{BEWEEG: D, (1, 3)} \\
            }
        }
    }
    & & \\ \\

    \theutterance \stepcounter{utterance}  
    & & & \multicolumn{2}{p{0.3\linewidth}}{
        \cellcolor[rgb]{0.9,0.9,0.9}{
            \makecell[{{p{\linewidth}}}]{
                \texttt{\tiny{[GM$|$GM]}}
                \texttt{BEWEEG: D, (1, 3)} \\
            }
        }
    }
    & & \\ \\

    \theutterance \stepcounter{utterance}  
    & & & \multicolumn{2}{p{0.3\linewidth}}{
        \cellcolor[rgb]{0.9,0.9,0.9}{
            \makecell[{{p{\linewidth}}}]{
                \texttt{\tiny{[GM$|$GM]}}
                \texttt{D moved to (1, 3)} \\
            }
        }
    }
    & & \\ \\

    \theutterance \stepcounter{utterance}  
    & & \multicolumn{4}{p{0.6\linewidth}}{
        \cellcolor[rgb]{0.9,0.9,0.9}{
            \makecell[{{p{\linewidth}}}]{
                \texttt{\tiny{[GM$\rangle$P2]}}
                \texttt{Verplaatst 'D' naar (1,3) met succes. Uw bijgewerkte raster ziet er als volgt uit:} \\
\texttt{```} \\
\texttt{ 1234567} \\
\texttt{╔═══════╗} \\
\texttt{║◌◌◌◌◌◌◌║ 1} \\
\texttt{║◌◌◌◌◌◌◌║ 2} \\
\texttt{║D◌P◌T┌─╢ 3} \\
\texttt{║◌OA┌─┤◌║ 4} \\
\texttt{║◌◌M│S├─╢ 5} \\
\texttt{╟───┼─┘◌║ 6} \\
\texttt{║◌◌◌│◌◌◌║ 7} \\
\texttt{╚═══╧═══╝} \\
\texttt{```} \\
\texttt{Je speelt momenteel ronde 15 van maximaal 28.} \\
\texttt{Jullie hebben gezamenlijk 18 van 16 straffen verzameld.} \\
\texttt{De andere speler heeft een voorwerp op zijn rooster verplaatst.} \\
\\ 
\\ 
\texttt{Wat is je volgende opdracht?} \\
            }
        }
    }
    & \\ \\

    \theutterance \stepcounter{utterance}  
    & & & \multicolumn{4}{p{0.6\linewidth}}{
        \cellcolor[rgb]{0.9,0.9,0.9}{
            \makecell[{{p{\linewidth}}}]{
                \texttt{\tiny{[GM$\langle$P2]}}
                \texttt{ZEG: Klaar?} \\
            }
        }
    }
     \\ \\

    \theutterance \stepcounter{utterance}  
    & & & \multicolumn{2}{p{0.3\linewidth}}{
        \cellcolor[rgb]{0.9,0.9,0.9}{
            \makecell[{{p{\linewidth}}}]{
                \texttt{\tiny{[GM$|$GM]}}
                \texttt{ZEG: Klaar?} \\
            }
        }
    }
    & & \\ \\

    \theutterance \stepcounter{utterance}  
    & & \multicolumn{4}{p{0.6\linewidth}}{
        \cellcolor[rgb]{0.9,0.9,0.9}{
            \makecell[{{p{\linewidth}}}]{
                \texttt{\tiny{[P1$\langle$GM]}}
                \texttt{Verplaatst 'D' naar (1,3) met succes. Uw bijgewerkte raster ziet er als volgt uit:} \\
\texttt{```} \\
\texttt{ 1234567} \\
\texttt{╔═══════╗} \\
\texttt{║◌◌◌◌◌◌◌║ 1} \\
\texttt{║◌◌◌◌◌◌◌║ 2} \\
\texttt{║D◌P◌T┌─╢ 3} \\
\texttt{║◌OA┌─┤S║ 4} \\
\texttt{║◌◌M│◌├─╢ 5} \\
\texttt{╟───┼─┘◌║ 6} \\
\texttt{║◌◌◌│◌◌◌║ 7} \\
\texttt{╚═══╧═══╝} \\
\texttt{```} \\
\texttt{Je speelt momenteel ronde 15 van maximaal 28.} \\
\texttt{Jullie hebben gezamenlijk 18 van 16 straffen verzameld.} \\
\texttt{De andere speler stuurde dit bericht:} \\
\\ 
\\ 
\texttt{"Klaar?"} \\
\\ 
\\ 
\texttt{Wat is je volgende opdracht?} \\
            }
        }
    }
    & \\ \\

    \theutterance \stepcounter{utterance}  
    & \multicolumn{4}{p{0.6\linewidth}}{
        \cellcolor[rgb]{0.9,0.9,0.9}{
            \makecell[{{p{\linewidth}}}]{
                \texttt{\tiny{[P1$\rangle$GM]}}
                \texttt{BEWEEG: S, (5, 5)} \\
            }
        }
    }
    & & \\ \\

    \theutterance \stepcounter{utterance}  
    & & & \multicolumn{2}{p{0.3\linewidth}}{
        \cellcolor[rgb]{0.9,0.9,0.9}{
            \makecell[{{p{\linewidth}}}]{
                \texttt{\tiny{[GM$|$GM]}}
                \texttt{BEWEEG: S, (5, 5)} \\
            }
        }
    }
    & & \\ \\

    \theutterance \stepcounter{utterance}  
    & & & \multicolumn{2}{p{0.3\linewidth}}{
        \cellcolor[rgb]{0.9,0.9,0.9}{
            \makecell[{{p{\linewidth}}}]{
                \texttt{\tiny{[GM$|$GM]}}
                \texttt{S moved to (5, 5)} \\
            }
        }
    }
    & & \\ \\

    \theutterance \stepcounter{utterance}  
    & & \multicolumn{4}{p{0.6\linewidth}}{
        \cellcolor[rgb]{0.9,0.9,0.9}{
            \makecell[{{p{\linewidth}}}]{
                \texttt{\tiny{[GM$\rangle$P2]}}
                \texttt{Je bericht is doorgestuurd naar de andere speler.} \\
\\ 
\\ 
\texttt{Je speelt momenteel ronde 16 van maximaal 28.} \\
\texttt{Jullie hebben gezamenlijk 18 van 16 straffen verzameld.} \\
\texttt{De andere speler heeft een voorwerp op zijn rooster verplaatst.} \\
\\ 
\\ 
\texttt{Wat is je volgende opdracht?} \\
            }
        }
    }
    & \\ \\

    \theutterance \stepcounter{utterance}  
    & & & \multicolumn{4}{p{0.6\linewidth}}{
        \cellcolor[rgb]{0.9,0.9,0.9}{
            \makecell[{{p{\linewidth}}}]{
                \texttt{\tiny{[GM$\langle$P2]}}
                \texttt{ZEG: Klaar?} \\
            }
        }
    }
     \\ \\

    \theutterance \stepcounter{utterance}  
    & & & \multicolumn{2}{p{0.3\linewidth}}{
        \cellcolor[rgb]{0.9,0.9,0.9}{
            \makecell[{{p{\linewidth}}}]{
                \texttt{\tiny{[GM$|$GM]}}
                \texttt{ZEG: Klaar?} \\
            }
        }
    }
    & & \\ \\

    \theutterance \stepcounter{utterance}  
    & & \multicolumn{4}{p{0.6\linewidth}}{
        \cellcolor[rgb]{0.9,0.9,0.9}{
            \makecell[{{p{\linewidth}}}]{
                \texttt{\tiny{[P1$\langle$GM]}}
                \texttt{Verplaatst 'S' naar (5,5) met succes. Uw bijgewerkte raster ziet er als volgt uit:} \\
\texttt{```} \\
\texttt{ 1234567} \\
\texttt{╔═══════╗} \\
\texttt{║◌◌◌◌◌◌◌║ 1} \\
\texttt{║◌◌◌◌◌◌◌║ 2} \\
\texttt{║D◌P◌T┌─╢ 3} \\
\texttt{║◌OA┌─┤◌║ 4} \\
\texttt{║◌◌M│S├─╢ 5} \\
\texttt{╟───┼─┘◌║ 6} \\
\texttt{║◌◌◌│◌◌◌║ 7} \\
\texttt{╚═══╧═══╝} \\
\texttt{```} \\
\texttt{Je speelt momenteel ronde 16 van maximaal 28.} \\
\texttt{Jullie hebben gezamenlijk 18 van 16 straffen verzameld.} \\
\texttt{De andere speler stuurde dit bericht:} \\
\\ 
\\ 
\texttt{"Klaar?"} \\
\\ 
\\ 
\texttt{Wat is je volgende opdracht?} \\
            }
        }
    }
    & \\ \\

    \theutterance \stepcounter{utterance}  
    & \multicolumn{4}{p{0.6\linewidth}}{
        \cellcolor[rgb]{0.9,0.9,0.9}{
            \makecell[{{p{\linewidth}}}]{
                \texttt{\tiny{[P1$\rangle$GM]}}
                \texttt{ZEG: klaar!} \\
            }
        }
    }
    & & \\ \\

    \theutterance \stepcounter{utterance}  
    & & & \multicolumn{2}{p{0.3\linewidth}}{
        \cellcolor[rgb]{0.9,0.9,0.9}{
            \makecell[{{p{\linewidth}}}]{
                \texttt{\tiny{[GM$|$GM]}}
                \texttt{ZEG: klaar!} \\
            }
        }
    }
    & & \\ \\

    \theutterance \stepcounter{utterance}  
    & & \multicolumn{4}{p{0.6\linewidth}}{
        \cellcolor[rgb]{0.9,0.9,0.9}{
            \makecell[{{p{\linewidth}}}]{
                \texttt{\tiny{[GM$\rangle$P2]}}
                \texttt{Je bericht is doorgestuurd naar de andere speler.} \\
\\ 
\\ 
\texttt{Je speelt momenteel ronde 17 van maximaal 28.} \\
\texttt{Jullie hebben gezamenlijk 18 van 16 straffen verzameld.} \\
\texttt{De andere speler stuurde dit bericht:} \\
\\ 
\\ 
\texttt{"klaar!"} \\
\\ 
\\ 
\texttt{Wat is je volgende opdracht?} \\
            }
        }
    }
    & \\ \\

    \theutterance \stepcounter{utterance}  
    & & & \multicolumn{4}{p{0.6\linewidth}}{
        \cellcolor[rgb]{0.9,0.9,0.9}{
            \makecell[{{p{\linewidth}}}]{
                \texttt{\tiny{[GM$\langle$P2]}}
                \texttt{ZEG: klaar!} \\
            }
        }
    }
     \\ \\

    \theutterance \stepcounter{utterance}  
    & & & \multicolumn{2}{p{0.3\linewidth}}{
        \cellcolor[rgb]{0.9,0.9,0.9}{
            \makecell[{{p{\linewidth}}}]{
                \texttt{\tiny{[GM$|$GM]}}
                \texttt{ZEG: klaar!} \\
            }
        }
    }
    & & \\ \\

    \theutterance \stepcounter{utterance}  
    & & & \multicolumn{2}{p{0.3\linewidth}}{
        \cellcolor[rgb]{0.9,0.9,0.9}{
            \makecell[{{p{\linewidth}}}]{
                \texttt{\tiny{[GM$|$GM]}}
                \texttt{true} \\
            }
        }
    }
    & & \\ \\

    \theutterance \stepcounter{utterance}  
    & & & \multicolumn{2}{p{0.3\linewidth}}{
        \cellcolor[rgb]{0.9,0.9,0.9}{
            \makecell[{{p{\linewidth}}}]{
                \texttt{\tiny{[GM$|$GM]}}
                \texttt{Initial states:} \\
\texttt{```} \\
\texttt{Player 1:} \\
\texttt{ 1234567} \\
\texttt{╔═══════╗} \\
\texttt{║◌◌◌◌◌T◌║ 1} \\
\texttt{║P◌◌◌◌◌◌║ 2} \\
\texttt{║◌◌AD◌┌─╢ 3} \\
\texttt{║◌◌◌┌─┤S║ 4} \\
\texttt{║O◌◌│◌├─╢ 5} \\
\texttt{╟───┼─┘◌║ 6} \\
\texttt{║◌M◌│◌◌◌║ 7} \\
\texttt{╚═══╧═══╝} \\
\\ 
\\ 
\texttt{Player 2:} \\
\texttt{ 1234567} \\
\texttt{╔═══════╗} \\
\texttt{║◌◌◌◌◌S◌║ 1} \\
\texttt{║◌◌◌◌O◌◌║ 2} \\
\texttt{║◌◌◌M◌┌─╢ 3} \\
\texttt{║◌◌A┌─┤◌║ 4} \\
\texttt{║◌◌◌│◌├─╢ 5} \\
\texttt{╟───┼─┘◌║ 6} \\
\texttt{║D◌T│P◌◌║ 7} \\
\texttt{╚═══╧═══╝} \\
\\ 
\texttt{```} \\
            }
        }
    }
    & & \\ \\

    \theutterance \stepcounter{utterance}  
    & & & \multicolumn{2}{p{0.3\linewidth}}{
        \cellcolor[rgb]{0.9,0.9,0.9}{
            \makecell[{{p{\linewidth}}}]{
                \texttt{\tiny{[GM$|$GM]}}
                \texttt{End states:} \\
\texttt{```} \\
\texttt{Player 1:} \\
\texttt{ 1234567} \\
\texttt{╔═══════╗} \\
\texttt{║◌◌◌◌◌◌◌║ 1} \\
\texttt{║◌◌◌◌◌◌◌║ 2} \\
\texttt{║D◌P◌T┌─╢ 3} \\
\texttt{║◌OA┌─┤◌║ 4} \\
\texttt{║◌◌M│S├─╢ 5} \\
\texttt{╟───┼─┘◌║ 6} \\
\texttt{║◌◌◌│◌◌◌║ 7} \\
\texttt{╚═══╧═══╝} \\
\\ 
\\ 
\texttt{Player 2:} \\
\texttt{ 1234567} \\
\texttt{╔═══════╗} \\
\texttt{║◌◌◌◌◌◌◌║ 1} \\
\texttt{║◌◌◌◌◌◌◌║ 2} \\
\texttt{║D◌P◌T┌─╢ 3} \\
\texttt{║◌OA┌─┤◌║ 4} \\
\texttt{║◌◌M│S├─╢ 5} \\
\texttt{╟───┼─┘◌║ 6} \\
\texttt{║◌◌◌│◌◌◌║ 7} \\
\texttt{╚═══╧═══╝} \\
\\ 
\texttt{```} \\
            }
        }
    }
    & & \\ \\

    \theutterance \stepcounter{utterance}  
    & & & \multicolumn{2}{p{0.3\linewidth}}{
        \cellcolor[rgb]{0.9,0.9,0.9}{
            \makecell[{{p{\linewidth}}}]{
                \texttt{\tiny{[GM$|$GM]}}
                \texttt{* success: True} \\
\texttt{* lose: False} \\
\texttt{* aborted: False} \\
\texttt{{-}{-}{-}{-}{-}{-}{-}} \\
\texttt{* Shifts: 7.00} \\
\texttt{* Max Shifts: 12.00} \\
\texttt{* Min Shifts: 6.00} \\
\texttt{* End Distance Sum: 0.00} \\
\texttt{* Init Distance Sum: 31.75} \\
\texttt{* Expected Distance Sum: 29.33} \\
\texttt{* Penalties: 18.00} \\
\texttt{* Max Penalties: 16.00} \\
\texttt{* Rounds: 17.00} \\
\texttt{* Max Rounds: 28.00} \\
\texttt{* Object Count: 7.00} \\
            }
        }
    }
    & & \\ \\

\end{supertabular}
}

\end{document}
